\documentclass[11pt]{article}

\usepackage[margin=1in]{geometry}
\usepackage{amsmath,amssymb}
\usepackage{graphicx}
\usepackage{booktabs}
\usepackage{caption}
\usepackage{float}
\usepackage{xcolor}
\usepackage{listings}
\usepackage[colorlinks=true,linkcolor=blue,citecolor=blue,urlcolor=blue]{hyperref}

\lstset{
  basicstyle=\ttfamily\footnotesize,
  columns=fullflexible,
  breaklines=true,
  frame=single,
  language=Python,
  showstringspaces=false,
  keywordstyle=\color{blue!70!black},
  commentstyle=\color{gray},
  stringstyle=\color{orange!70!black}
}

\title{A Leakage-Safe Evaluation Pipeline for Heart/Lung Sound\\
Source Separation on HLS-CMDS}
\author{[Author Name] \\ TEEP Internship}
\date{2026-08-17, revised 2026-09-13}

\begin{document}
\maketitle

\begin{abstract}
This report documents a software pipeline built to evaluate blind source
separation of heart and lung sounds on the HLS-CMDS dataset
\cite{hlscmds}. The pipeline consists of: (1) a data-loading and
integrity-validation layer, including a per-row additive-triplet audit that
found only 36 of \texttt{Mix.csv}'s 145 rows have a mixed recording that is
acoustically related to its named heart/lung sources at all --- the other
109 fail even before leakage is considered; (2) a leakage-safe,
triplet-level cross-validation split derived from a content-hash audit of
the dataset that found a third of the ``isolated'' ground-truth recordings
are literal duplicates of recordings used inside the mixtures; (3) a BSS
Eval (SDR/SIR/SAR) metrics module wrapping \texttt{mir\_eval}, validated
both by direct cross-checking against \texttt{mir\_eval} and by closed-form
synthetic tests; (4) a k-fold evaluation harness tying the split and the
metrics together with per-fold, across-fold, and pooled row-level
aggregation; (5) a combinatorial synthetic mixture set
(\texttt{synthetic\_mix.py}) that supersedes the 36-row native-additive set
as the primary evaluation substrate for a continuous SDR-vs-difficulty
sweep, with its own source-file-level leakage-safe split --- the
triplet/leak-group split in (2) provably collapses to a single connected
component once every heart file is combinatorially paired with every lung
file, so the split unit for this substrate had to move from the mixture row
up to the individual recording; and (6) six separation baselines run
end-to-end through the harness --- fixed bandpass filtering, a leakage-safe
reproduction of a published supervised-NMF method, an unsupervised
per-mixture NMF ablation, multi-stage Singular Spectrum Analysis (MSSA), a
reproduction of a published EVMD-based method, and Conv-TasNet-lite, the
first neural, end-to-end trained baseline. We report the audio-quality
characteristics of the dataset, the structure of the recording-reuse and
mixture-validity problems the split and the additive-triplet audit were
built to solve, and all six baselines' 5-fold SDR/SIR/SAR results on the
full native row set, the 36-row native-valid subset, and the 1500-row
synthetic substrate, and situate this work against six related papers
reviewed, none of which combine separation-quality reporting, isolated-vs.\
separated classification comparison, and a leakage-safe split on this
dataset.

Building on that separation-quality pipeline, this revision adds (7) a
downstream heart-sound classifier (MFCC features, an SVM), evaluated on
the identical leak-group folds, with a class-grouping decision that
reproduces Yaqub et al.'s own four-class scheme so results are directly
comparable to their motivating 89\%$\to$41\% result; (8) a controlled
degradation scheme that interpolates continuously between a separation
method's real output and the isolated ground truth, letting classification
accuracy be measured as a continuous function of separation quality (SDR)
rather than at two fixed endpoints, together with the generated sweep this
produces across all six baselines; (9) a weight-sharing protocol
(``Condition B'') for classifying separated audio with the exact same
trained classifier used on isolated audio, and the resulting
accuracy-vs-SDR curve and knee-point definition --- the SDR below which
separation stops beating not separating at all; (10) a second,
architecturally distinct classifier (a log-Mel spectrogram + CNN) used as
a robustness check on whether a knee point is a property of separation
quality or an artifact of one architecture; and (11) a computational-cost
analysis (multiply-accumulate operation counts and measured desktop
latency under a uniform protocol) placing every baseline's separation
quality against its compute cost on a Pareto plane. The six-baseline
sweep, the resulting knee-point curve and the second architecture's own
curve are now complete and reported in full (Sections~\ref{sec:conditionb}--\ref{sec:robustness});
an earlier revision reported them as partial results, and the
oracle-mask ceiling check that revision listed as informal is now a
committed, re-runnable script (Section~\ref{sec:discussion}).

\textbf{The dataset's provenance was verified independently after this
pipeline was built (Section~\ref{sec:provenance}): the working copy is
byte-identical to the Mendeley release, so the 4000\,Hz sample rate and the
109/145 mixture rows that fail the named-reference additivity test are
properties of HLS-CMDS as published rather than of a downsampled redistribution.}
\end{abstract}

\tableofcontents

% ============================================================
\section{Introduction}
% ============================================================

Heart and lung sounds are conventionally recorded and classified in
isolation, but a clinical stethoscope in practice picks up both
simultaneously. HLS-CMDS \cite{hlscmds} is, to our knowledge, the first
published dataset to offer matched isolated and mixed cardiorespiratory
recordings from the same underlying sources, which makes it possible to ask
a question the isolated-only or mixed-only datasets used elsewhere in the
literature cannot answer directly: \emph{how much does separating heart from
lung sound, before classifying it as normal or abnormal, cost in accuracy
relative to classifying the isolated ground truth?}

That question is not rhetorical. Yaqub et al.\ \cite{spectrotemporal}
measured exactly this cost on this project's own dataset and found it
severe: a PCG
classifier retrained on HLS-CMDS reaches 89.0\% accuracy on clean, isolated
recordings, but collapses to 41.0\% when evaluated on heart sounds
computationally separated from HLS-CMDS's mixed recordings with a standard
bandpass filter (their Experiment 4; see Section~\ref{sec:yaqub}). That
89\%$\to$41\% gap is the motivating result for this whole project: this
report's separation-quality metrics (SDR/SIR/SAR, Section~\ref{sec:metrics})
are what let that gap be explained mechanistically and traced out
continuously against separation quality, rather than measured at only the
two endpoints Yaqub et al.\ report.

Section~\ref{sec:background} gives a from-scratch, student-level
introduction to every technical idea the rest of this report depends on
(blind source separation, BSS Eval, NMF, SSA, VMD, Conv-TasNet's
convolutional architecture, cross-validation and data leakage,
classification metrics, and compute-cost measurement) so later sections
can explain \emph{why} a specific choice was made without re-deriving the
underlying concept every time. Section~\ref{sec:dataset} describes the
dataset, the audio-quality audit run over it, and a provenance problem
discovered after the pipeline below was built. Section~\ref{sec:leakage}
presents the recording-reuse finding that motivates a non-trivial
train/evaluation split. Section~\ref{sec:validity} presents a second,
independent data problem: most mixture rows turn out not to be
acoustically related to their named sources at all.
Section~\ref{sec:method} describes the software components built to
support this line of research: the split, the separation-quality metrics
and their validation, the cross-validation harness, and all six separation
baselines. Section~\ref{sec:results} reports results on the full row set,
the 36-row valid subset, and the 1500-row synthetic substrate.
Section~\ref{sec:classifier} builds on that separation-quality pipeline
with a downstream classifier, the controlled degradation scheme and SDR
sweep this project's own C2 research question needs, the resulting
accuracy-vs-SDR curve and knee point, and a second-architecture robustness
check on that knee point. Section~\ref{sec:compute} measures every
method's computational cost (MACs and desktop latency) and places
separation quality against compute cost on one plane.
Section~\ref{sec:related} surveys related work. Section~\ref{sec:discussion}
discusses the results and their implications, and
Section~\ref{sec:conclusion} summarizes open work.

All software described here lives in a single repository, organized as:
\texttt{src/load\_dataset.py} (loading/validation, including the
additive-triplet audit of Section~\ref{sec:validity}), \texttt{src/split.py}
(leakage-safe folds), \texttt{src/metrics.py} (BSS Eval), \texttt{src/}
\texttt{eval\_harness.py} (cross-validation), \texttt{src/baselines.py}
(both separation baselines), \texttt{src/statistics/audio\_quality.py}
(dataset audit), and \texttt{src/visualization/} (plotting scripts), with
\texttt{src/test\_metrics.py}, \texttt{src/test\_split.py}, and
\texttt{src/test\_load\_dataset.py} as the accompanying test suites (31
tests total, all passing at the time of writing).

% ============================================================
\section{Background and Preliminaries}
\label{sec:background}
% ============================================================

This section is written for a reader who knows general programming and
basic statistics but has not necessarily studied audio signal processing,
matrix factorization, or machine learning evaluation methodology before.
Every later section builds on one of the ideas explained here, and refers
back to this section by name rather than re-deriving the idea inline. If a
later section says ``because SAR captures this'' or ``this would be data
leakage,'' the explanation of \emph{why} that matters is here, not
repeated each time.

\subsection{The problem: blind source separation}
\label{sec:bg-bss}

A stethoscope placed on a chest picks up heart sounds and lung sounds at
the same time, physically summed into one waveform by the air and tissue
they travel through before reaching the microphone. Call the true heart
waveform $h(t)$, the true lung waveform $l(t)$, and the recorded mixture
$m(t)$. The simplest physical model of what happened is a linear mix:
\begin{equation}
m(t) = a\cdot\bigl(h(t) + l(t)\bigr) + n(t),
\label{eq:linearmix}
\end{equation}
where $a$ is an overall recording gain and $n(t)$ is any additional noise.
\emph{Source separation} means recovering estimates $\hat h(t)$ and
$\hat l(t)$ from $m(t)$ alone. It is called \emph{blind} when the
separation method is not told what $h$ and $l$ individually sound like for
this specific recording --- it only gets $m(t)$, plus at most some general
knowledge about what heart sounds and lung sounds tend to look like
(learned from \emph{other}, separate recordings).

This is a fundamentally underdetermined problem: Equation~\ref{eq:linearmix}
is one equation with two unknown functions on the right-hand side, so there
are infinitely many pairs $(\hat h, \hat l)$ that add up to exactly the
observed $m$. The only reason separation is possible at all is that heart
and lung sounds are not arbitrary --- they differ systematically in
\emph{where their energy sits}, whether that is described in frequency
(Sections~\ref{sec:baseline1}, \ref{sec:baseline4}, \ref{sec:baseline5}), in
a learned spectral ``dictionary'' (Sections~\ref{sec:baseline2},
\ref{sec:baseline3}), or in patterns a neural network discovers on its own
from examples (Section~\ref{sec:baseline6}). Every separation method in
this report is, underneath its specific mathematics, a different way of
exploiting that systematic difference to break the underdetermined problem.

\subsection{How do you know if a separation worked? BSS Eval (SDR/SIR/SAR)}
\label{sec:bg-bsseval}

Once a method produces an estimate $\hat s_j$ for source $j$ (heart or
lung), the natural first idea is to compare it to the true source $s_j$ with
something like mean-squared error. That turns out to be a poor measure for
audio, because it does not distinguish \emph{why} $\hat s_j$ differs from
$s_j$ --- a barely-audible time shift and a completely wrong signal can
produce the same squared error. The BSS Eval framework of Vincent et al.\
\cite{vincent2006} (Section~\ref{sec:relatedmetrics}) fixes this by
decomposing the error into three parts with different \emph{causes}:
\begin{equation}
\hat s_j = s_{\text{target}} + e_{\text{interf}} + e_{\text{artif}},
\end{equation}
where $s_{\text{target}}$ is the best time-invariant-filtered version of
the \emph{true} source $s_j$ that could be extracted from the estimate
(so a harmless gain change or a short, fixed delay does not get punished),
$e_{\text{interf}}$ is the part of the estimate's error that is explained
by \emph{other sources leaking in} (e.g.\ lung energy showing up in the
heart estimate), and $e_{\text{artif}}$ is whatever is left over ---
distortion the separation algorithm itself introduced that cannot be
explained by leftover interference at all (filter ringing, NMF
reconstruction error, neural-network hallucination). Three ratios follow
directly, each answering a different question, always expressed in
decibels (a base-10 logarithmic scale, $10\log_{10}(\cdot)$, chosen because
audio energy ratios span many orders of magnitude and human loudness
perception is itself roughly logarithmic --- a linear ratio would make
large, easy-to-separate cases and small, hard cases look deceptively
similar on a plot):
\begin{align}
\text{SDR} &= 10\log_{10}\frac{\lVert s_{\text{target}} \rVert^2}
  {\lVert e_{\text{interf}} + e_{\text{artif}} \rVert^2}
  &\text{``how good is the estimate, overall?''} \\
\text{SIR} &= 10\log_{10}\frac{\lVert s_{\text{target}} \rVert^2}
  {\lVert e_{\text{interf}} \rVert^2}
  &\text{``how much of the other source leaked in?''} \\
\text{SAR} &= 10\log_{10}\frac{\lVert s_{\text{target}} + e_{\text{interf}} \rVert^2}
  {\lVert e_{\text{artif}} \rVert^2}
  &\text{``how much did the method itself distort the signal?''}
\end{align}
\textbf{Why report all three instead of just SDR.} A method can score well
on one and poorly on another, and the gap is itself diagnostic. An
aggressive bandpass filter, for instance, blocks the other source
effectively (good SIR) but also rings and destroys the target's own
harmonics in the process (poor SAR); a method that barely does anything
scores the opposite way. SDR folds both effects into one number, which
makes it convenient but also means a good SDR can hide a bad SAR that a
downstream classifier might actually care about (Section~\ref{sec:relatedmetrics}'s
point from Vincent et al.\ that a classifier can stay accurate despite real
artifacts, or fail because of them, depending on whether the artifact
crosses a decision boundary) --- reporting only SDR would erase exactly the
kind of trade-off Section~\ref{sec:results} needs to describe.

\subsection{Matrix factorization for audio: Non-negative Matrix Factorization (NMF)}
\label{sec:bg-nmf}

Baselines~2 and~3 (Sections~\ref{sec:baseline2}--\ref{sec:baseline3}) work
not on the raw waveform but on its \emph{spectrogram}: split the waveform
into short overlapping windows (a Short-Time Fourier Transform, STFT), take
the magnitude of each window's frequency spectrum, and stack these columns
side by side into a matrix $V$ of shape (frequency bins) $\times$ (time
frames). Every entry of $V$ is a non-negative energy value --- there is no
such thing as negative sound energy at a given frequency and time.

Non-negative Matrix Factorization asks: can $V$ (size $F\times T$) be
approximated as the product of two much smaller non-negative matrices,
$V \approx WH$, where $W$ (size $F\times K$) holds $K$ ``basis spectra''
(recurring frequency shapes) and $H$ (size $K\times T$) holds, for every
time frame, how strongly each basis spectrum is active? $K$ is chosen much
smaller than $F$ or $T$, so this is a compression that forces the
algorithm to explain the whole spectrogram using only $K$ reusable
building blocks. The non-negativity constraint is what makes the result
interpretable as a \emph{parts-based} representation (unlike, say, PCA/SVD
factorizations, which allow negative combinations and produce components
that only make sense added and subtracted together, not individually):
every basis spectrum in $W$ is something a real sound could actually look
like, and every column of $H$ says ``how much of each real sound is
present right now,'' non-negatively.

Fitting $W$ and $H$ is done by \emph{multiplicative updates} (Lee \&
Seung, 2001): starting from a random non-negative guess, $W$ and $H$ are
repeatedly rescaled by a ratio that is guaranteed to be non-negative and
that provably decreases the reconstruction error (specifically the
KL-divergence between $V$ and $WH$ used here) at every step. This is a
form of coordinate descent --- it is simple to implement and, unlike
generic gradient descent, never needs a learning rate or risks producing a
negative value that would have to be clipped.

\textbf{Supervised NMF} (Baseline~2, Section~\ref{sec:baseline2}) exploits
this by learning $W$ \emph{separately} for heart and lung, from clean,
isolated recordings of each, before ever looking at a mixture --- exactly
analogous to learning a small palette of ``basis colors'' from pure paint
samples. The two $W$ matrices are then frozen and stacked together, and
only $H$ is re-solved (with the same multiplicative-update rule, holding
$W$ fixed) against a specific held-out \emph{mixture}'s spectrogram. This
answers ``how much of each already-known basis spectrum would need to be
present, at each moment, to reconstruct this new, unlabeled mixture?'' ---
the same idea as decomposing an unknown mixed color as a combination of
known primary colors. Once $H$ says how much of the heart-dictionary and
lung-dictionary explain the mixture at each time-frequency point, that
proportion becomes a \emph{soft mask}: the fraction of the mixture's
complex spectrogram at each point assigned to heart is
$W_{\text{heart}}H_{\text{heart}} /
(W_{\text{heart}}H_{\text{heart}}+W_{\text{lung}}H_{\text{lung}})$, applied
directly to the mixture's own (complex-valued) STFT before inverting back
to a waveform --- this is a \emph{Wiener-style} mask, the standard way to
turn an energy decomposition into an actual audio estimate. \textbf{Standard
NMF} (Baseline~3, Section~\ref{sec:baseline3}) is the same machinery with
one piece removed: no pre-learned dictionary at all, both $W$ and $H$ are
factorized fresh out of each mixture's own spectrogram (a purely
\emph{blind}, or \emph{transductive}, use of NMF), which isolates how much
of Baseline~2's performance is actually attributable to its learned,
frozen dictionary rather than to NMF factorization itself.

\subsection{An alternative decomposition: Singular Spectrum Analysis (SSA)}
\label{sec:bg-ssa}

Baseline~4 (Section~\ref{sec:baseline4}) uses a completely different
mathematical tool: Singular Spectrum Analysis. The construction has three
steps. First, \emph{embedding}: turn the 1-D waveform of length $N$ into a
2-D matrix by sliding a window of length $L$ along it one sample at a time
and stacking each window as a column, giving an $L\times(N-L+1)$
\emph{trajectory matrix}. This turns ``what happens over time'' into ``what
does this matrix's structure look like,'' which linear algebra has strong
tools for. Second, \emph{decomposition}: apply the Singular Value
Decomposition (SVD), which writes any matrix as a sum of rank-1 pieces
(each an outer product of one ``time pattern'' vector and one ``lag
pattern'' vector) ordered by how much of the trajectory matrix's total
energy each piece explains. These rank-1 pieces are the same mathematical
object behind PCA, and each one corresponds to some characteristic
oscillation frequency in the original signal --- SSA does not force this
frequency interpretation, but it falls out naturally because a matrix
built from a windowed periodic signal has exactly this kind of low-rank
structure. Third, \emph{diagonal averaging} (Hankelization): each selected
rank-1 piece, still shaped like an $L\times(N-L+1)$ matrix, is converted
back into a 1-D \emph{reconstructed component} (RC) by averaging along its
anti-diagonals, undoing the windowing from step one.

Given a set of RCs, each with its own dominant frequency, separating heart
from lung reduces to sorting them by that frequency and grouping the
low-frequency ones into a cardiac output and the higher-frequency ones into
a respiratory output --- which is exactly what MSSA's two-stage design does
(a first pass with a fixed cardiac/respiratory frequency threshold, then a
second SSA pass on the leftover residual to more carefully pick out the
respiratory content by how much variance each RC explains and how
correlated it is with what has already been selected). No dictionary is
learned from other recordings anywhere in this process --- SSA operates
entirely on the one signal it is given, which is why Baseline~4 (like
Baseline~1) is a zero-training method with no leakage exposure at all
(Section~\ref{sec:leakage}).

\subsection{Optimization-based decomposition: Variational Mode Decomposition (VMD)}
\label{sec:bg-vmd}

Baseline~5 (Section~\ref{sec:baseline5}) uses Variational Mode
Decomposition (Dragomiretskiy \& Zosso, 2014), which frames separation as
an explicit optimization problem rather than a matrix factorization or a
fixed filter: find $K$ ``modes,'' each a signal that is as narrow-band as
possible (concentrated tightly around its own, simultaneously-discovered
center frequency $\omega_k$), subject to the constraint that all $K$ modes
sum back to the original signal exactly. ``As narrow-band as possible'' is
made precise by minimizing the modes' total bandwidth (measured via the
signal's Hilbert transform, which shifts a real signal's spectrum so its
energy content at each frequency can be isolated and penalized). This
joint problem --- $K$ unknown signals \emph{and} $K$ unknown center
frequencies, all coupled through one shared constraint --- is solved by
the Alternating Direction Method of Multipliers (ADMM), a general
optimization strategy for exactly this situation: instead of solving for
everything at once, alternate between updating each mode (holding all
others and the frequencies fixed) and updating each center frequency
(holding the modes fixed), repeating until the modes stop changing. Once
the mode with the lowest center frequency is identified as the cardiac
component, subtracting it from the mixture gives the respiratory estimate.
$K$ itself is not known in advance and is swept from 2 upward until a
paper-specified quality check passes (or a ceiling is hit) --- a further
practical detail documented in Section~\ref{sec:baseline5}.

\subsection{Deep learning for separation: convolutions, dilation, and Conv-TasNet}
\label{sec:bg-convtasnet}

Baseline~6 (Section~\ref{sec:baseline6}) replaces a hand-designed
transform (STFT, trajectory matrix, mode decomposition) with weights
learned end-to-end from data. Its core operation is the 1-D convolution: a
short, learned filter (a small vector of weights) is slid across the
waveform, and at each position it computes a weighted sum of the samples
underneath it --- conceptually the same sliding-window idea as
Baseline~1's Butterworth filter (Section~\ref{sec:baseline1}), except the
filter's coefficients are \emph{learned} from examples rather than
hand-derived from a PSD survey. Stacking many such convolutional layers,
each one operating on the previous layer's output, lets the network build
up increasingly abstract, longer-range patterns from short, local ones.

A plain stack of convolutions would need very many layers to ``see'' a
useful amount of context (each layer only extends what the next one can
see by roughly its own filter length), which is where \emph{dilated}
convolutions come in: a dilated filter skips input positions at a fixed
spacing, so the same small number of weights covers a much wider span of
input without adding parameters. Stacking layers with exponentially
increasing dilation (as Conv-TasNet's Temporal Convolutional Network, TCN,
does: dilation $2^0, 2^1, 2^2,\dots$) lets the network's effective
\emph{receptive field} --- how much of the input any one output value
depends on --- grow exponentially with depth instead of linearly, which is
how Baseline~6 covers roughly 1.27\,s of audio context with only a
handful of layers (Section~\ref{sec:baseline6}).

Conv-TasNet's overall architecture has three pieces trained jointly:
an \emph{encoder} (a bank of learned 1-D convolution filters that plays the
same role the STFT plays for the NMF baselines, except the filter shapes
are learned rather than fixed sinusoids), a \emph{separator} (the dilated
TCN described above, which looks at the encoded mixture and predicts one
soft mask per source), and a \emph{decoder} (the mathematical inverse of
the encoder, turning a masked encoded representation back into a
waveform). Training this end-to-end means the network never has to be told
what a ``basis spectrum'' or a ``frequency band'' is --- it discovers
whatever representation and masking strategy minimizes its training loss
directly, which is both Baseline~6's main advantage (it can in principle
discover structure the other five methods' fixed transforms cannot
represent) and its main practical cost (it needs enough labeled training
examples to find that structure reliably, which Section~\ref{sec:baseline6}
and Section~\ref{sec:results} both flag as a real limitation on this
dataset's small recording pool).

\subsection{Evaluating machine learning fairly: cross-validation and data leakage}
\label{sec:bg-cv}

Any time a method is \emph{fit} to data (a learned dictionary, a trained
network, even a hand-tuned filter cutoff), evaluating it on the same data
it was fit on overstates how well it will do on genuinely new recordings
--- the method may simply have memorized quirks of that specific data
rather than learned something general. The standard fix is to hold out
some data, fit only on the rest, and evaluate only on the held-out part.
\emph{$k$-fold cross-validation} extends this so that every data point
gets to be in the held-out set exactly once (across $k$ different
``folds''), so the whole dataset is used for both fitting and evaluation
without ever evaluating a fold's result using data it was fit on ---
important on a dataset this small, where reserving a large, fixed holdout
would leave too little data to fit anything meaningful.

\emph{Data leakage} is the general name for evaluating a model in a way
that lets information from the held-out set influence what it was fit on,
even indirectly, which silently inflates the measured score without the
model actually generalizing any better. The textbook example is fitting
and evaluating on data that turn out to contain duplicate rows: if a
near-copy of a test example sits in the training set, the model does not
have to generalize to score well on it. This project's leak-group split
(Section~\ref{sec:leakage}) exists precisely because HLS-CMDS has exactly
this problem in a non-obvious form: the same physical recording appears
multiple times under \emph{different ID strings} (once standalone in
\texttt{HS.csv}/\texttt{LS.csv}, and again inside one or more
\texttt{Mix.csv} rows), so a standard, generic $k$-fold split --- which
only knows about row IDs, not about which underlying audio content those
IDs point to --- would silently leak a mixture's own ground truth into the
pool used to fit its separation dictionary. Detecting this required
hashing raw audio content directly rather than trusting the ID columns
(Section~\ref{sec:leakage}), and the resulting split unit is not the row
but a \emph{leak group} (a connected component of rows that share
audio) --- the minimal change to standard cross-validation that actually
closes this specific leak.

\subsection{Classification metrics: accuracy, macro-F1, and the confusion matrix}
\label{sec:bg-classmetrics}

Section~\ref{sec:classifier} onward evaluates a \emph{classifier} (does a
heart recording belong to Normal, Murmur, Extra Sound, or Rhythm
Disorder?), which needs its own vocabulary. \emph{Accuracy} is simply the
fraction of examples classified correctly, but it has a well-known blind
spot on imbalanced data: if 24 of 50 recordings are Murmur and a
classifier always predicts ``Murmur'' regardless of the input, it scores
48\% accuracy while having learned nothing useful about the other three
classes. \emph{Macro-F1} is one standard fix: compute the F1 score
(the harmonic mean of precision and recall) separately for \emph{each}
class, then average those per-class scores with equal weight --- so a
model that nails the common class but fails every rare one no longer gets
to hide behind the common class's size. A \emph{confusion matrix} is the
raw material both metrics are computed from: a table whose rows are the
true class and columns are the predicted class, so cell $(i,j)$ counts how
many examples of true class $i$ were predicted as class $j$. The diagonal
is what a classifier got right; everything off the diagonal is a specific,
nameable mistake (``Rhythm Disorder recordings are most often predicted as
Murmur''), which is far more actionable than a single aggregate number and
is exactly the level of detail Section~\ref{sec:classifier-failure} uses
to describe each architecture's specific failure pattern.

\subsection{Computational cost: MACs, latency, and the Pareto front}
\label{sec:bg-compute}

Section~\ref{sec:compute} asks a different kind of question: not ``how
good is this method,'' but ``how expensive is it to run.'' The atomic unit
of cost used throughout that section is the \emph{multiply-accumulate}
(MAC): one multiplication followed by one addition into a running total,
which is the single operation almost every method in this report is built
out of at its core, whether it is a matrix multiplication inside NMF, a
convolution inside Conv-TasNet, or an SVD inside SSA. Counting MACs gives a
hardware-independent proxy for ``how much arithmetic work'' a method
requires --- useful because it can be computed analytically from a
method's own equations, without needing to actually run it. It is,
however, only a \emph{proxy}: how quickly a given MAC count actually
executes on real hardware depends heavily on how well that arithmetic maps
onto the hardware's own optimized code paths (vectorized instructions,
cache-friendly memory access, a mature compiled library vs.\ a direct,
unoptimized implementation of a paper's own pseudocode). Section~\ref{sec:compute}
measures both MACs (the theoretical proxy) and real wall-clock latency
(the measured outcome) \emph{separately}, specifically because they turn
out to disagree for at least one method here --- a divergence that a MACs
count alone could never reveal.

Comparing several methods on two axes at once (quality and cost) uses the
idea of a \emph{Pareto front}, borrowed from multi-objective optimization:
method A is said to \emph{dominate} method B if A is at least as good as B
on every axis being compared and strictly better on at least one. A
dominated method is never the right choice, no matter how the two axes are
weighted against each other, because some other method is simply better in
every way that matters; a method that is \emph{not} dominated by anything
represents a genuine trade-off (better quality at higher cost, or lower
cost at some quality loss) that a deployment decision would actually have
to weigh. Section~\ref{sec:compute-plane} applies this directly to the six
separation baselines' SDR-vs-compute trade-off.

% ============================================================
\section{Dataset}
\label{sec:dataset}
% ============================================================

HLS-CMDS \cite{hlscmds} records heart and lung sounds from a CAE Juno
mid-fidelity patient simulator with swappable male/female chest skins,
captured with a digital stethoscope in a quiet room. Heart sounds are
recorded at six chest landmarks (Apex, RUSB, LUSB, LLSB, RC, LC) across ten
sound-type classes (Normal, Late Diastolic Murmur, Mid Systolic Murmur, Late
Systolic Murmur, Atrial Fibrillation, S4, Early Systolic Murmur, S3,
Tachycardia, AV Block); lung sounds are recorded at six anterior landmarks
(LLA, LMA, LUA, RLA, RMA, RUA) across six classes (Normal, Wheezing, Fine
Crackles, Rhonchi, Pleural Rub, Coarse Crackles). The dataset ships three
CSV manifests: \texttt{HS.csv} (50 isolated heart recordings),
\texttt{LS.csv} (50 isolated lung recordings), and \texttt{Mix.csv} (145
mixed recordings, each row also naming the heart and lung recording used to
build that mixture).

\subsection{Dataset provenance: verified against the release}
\label{sec:provenance}

\textbf{The dataset copy every result in this report is computed on is the
released one.} The working copy in this project (\texttt{HLS\_CMDS/}) was
obtained from
\href{https://github.com/Torabiy/HLS-CMDS}{github.com/Torabiy/HLS-CMDS}
rather than the Mendeley release
(\href{https://data.mendeley.com/datasets/8972jxbpmp/3}{data.mendeley.com/datasets/8972jxbpmp/3})
the project charter specifies, and an internal review
(\emph{TEEP2026\_Sprint0\_Review}, 2026-08-18) flagged the resulting
possibility that everything here had been measured on a downsampled
convenience mirror. That possibility was tested directly and did not hold: a
fresh Mendeley download (v3, retrieved 2026-08-19), checksummed independently
of this project's own tooling, is byte-identical to the working copy on all
six top-level artifacts (Table~\ref{tab:checksums}).

\begin{table}[H]
\centering
\caption{Provenance check. Mendeley v3, retrieved 2026-08-19; MD5 over each
top-level artifact, identical between the fresh download and the working
copy.}
\label{tab:checksums}
\begin{tabular}{@{}ll@{}}
\toprule
Artifact & MD5 \\
\midrule
\texttt{HS.csv} & \texttt{b61a866d2e8d70ae2910e4c049327d1e} \\
\texttt{HS.zip} & \texttt{7d678aff168184c13d8893b3bef4f4ba} \\
\texttt{LS.csv} & \texttt{fa40494d1dcbcf173ee5e9583b3faf82} \\
\texttt{LS.zip} & \texttt{d56a0d8ce5941342f33bfa9b04ef2bd9} \\
\texttt{Mix.csv} & \texttt{377b55ea6361ff7c48994cddaf5c6e6f} \\
\texttt{Mix.zip} & \texttt{f28243aa5394324c58f90a466d60c2f4} \\
\bottomrule
\end{tabular}
\end{table}

Two findings survive that check as properties of the published dataset, which
is why this subsection is retained rather than deleted:

\begin{itemize}
  \item \textbf{Sample rate.} Every file is 4000\,Hz
    (Section~\ref{sec:audioquality}), not the 22{,}050\,Hz the HLS-CMDS
    descriptor paper states --- a $5.5\times$ mismatch belonging to the
    release itself. Whether the descriptor misstates the rate or the figure
    was mis-transcribed downstream of it is not resolved here. It is not a
    problem for the science, since heart energy sits at 20--200\,Hz and lung
    at 100--1000\,Hz (Section~\ref{sec:baseline1}), comfortably under a
    2000\,Hz Nyquist; it is a problem for tooling, and
    Section~\ref{sec:baseline6} documents the consequence.
  \item \textbf{Mixture validity.} 109 of 145 \texttt{Mix.csv} rows'
    \emph{mixed} recording is acoustically unrelated to its named heart and
    lung sources (Section~\ref{sec:validity}); only 36 rows are usable as
    paired separation ground truth. Re-downloading cannot recover them,
    because the re-download is the same bytes.
\end{itemize}

Both figures come from HS/LS/Mix.csv's own audio content, not from metadata,
so they are not artifacts of a mislabeled manifest. Every
separation-quality number in this report is accordingly a measurement on
HLS-CMDS as published.

\subsection{Audio-quality audit}
\label{sec:audioquality}

\texttt{src/statistics/audio\_quality.py} scans every \texttt{.wav} file
referenced by the three manifests (535 files total: 50 + 50 + 145
$\times$ 3, since each mix row's heart, lung, and mixed components are
separate files) and reports duration, sample rate, channel count, and
clipping (samples pinned at the 16-bit PCM full-scale rail,
$\pm 32767$), aggregated per sound-type class and per recording location.

\textbf{Format.} Every file is 4000\,Hz, mono, exactly 15.0\,s
(60{,}000 samples) --- perfectly uniform, no resampling or trimming
required before feature extraction. This uniformity is itself part of the
provenance discrepancy: 4000\,Hz is not the sample rate the HLS-CMDS
descriptor paper documents (Section~\ref{sec:provenance}).

\textbf{Clipping.} 36 of 535 files (6.7\%) clip, but only by a single sample
each ($1/60{,}000 \approx 0.0017\%$ of the file), and exclusively in the
\emph{mixed} component of \texttt{Mix.csv} rows (never in \texttt{HS.csv},
\texttt{LS.csv}, or the mix set's own heart/lung components) --- specifically
row M0087 and the contiguous run M0111--M0145. This set of 36 rows is exactly
the set Section~\ref{sec:validity}'s additive-triplet audit independently
finds acoustically valid, which resolves the mechanism: these are the rows
where \texttt{mixed} was actually built by summing the heart and lung
waveforms and then peak-normalizing the result (an observed gain of
3.5--135$\times$ across the 36 rows), and peak normalization of an
already-full-scale sum pins one sample at the 16-bit rail. It is a
normalization fingerprint, not a recording defect, and it has not been
reported upstream.

\subsubsection{Per-class and per-location breakdown}
\label{sec:perclass}

Beyond the dataset-wide summary above, \texttt{audio\_quality.py} also
aggregates duration, sample rate, and clipping \emph{per class} and \emph{per
recording location}, over the 100 isolated recordings (\texttt{HS.csv} +
\texttt{LS.csv}). Duration (15.0\,s, std 0.0) and sample rate (4000\,Hz) are
constant across every class and every location --- omitted from
Tables~\ref{tab:perclass}--\ref{tab:perlocation} below since there is no
variation to report --- so the only columns that carry information are the
class/location's sample count and its clipped-file count.

\begin{table}[H]
\centering
\caption{Per-class breakdown, isolated recordings only (\texttt{HS.csv}/\texttt{LS.csv}).}
\label{tab:perclass}
\begin{tabular}{@{}llc c@{}}
\toprule
Component & Class & $n$ & Clipped \\
\midrule
Heart & Normal                 & 9  & 0 \\
Heart & Mid Systolic Murmur    & 7  & 0 \\
Heart & Late Diastolic Murmur  & 6  & 0 \\
Heart & Early Systolic Murmur  & 6  & 0 \\
Heart & Late Systolic Murmur   & 5  & 0 \\
Heart & S3                     & 5  & 0 \\
Heart & Atrial Fibrillation    & 4  & 0 \\
Heart & AV Block               & 3  & 0 \\
Heart & Tachycardia            & 3  & 0 \\
Heart & S4                     & 2  & 0 \\
\addlinespace
Lung  & Normal            & 12 & 0 \\
Lung  & Coarse Crackles   & 9  & 0 \\
Lung  & Pleural Rub       & 9  & 0 \\
Lung  & Rhonchi           & 8  & 0 \\
Lung  & Wheezing          & 7  & 0 \\
Lung  & Fine Crackles     & 5  & 0 \\
\bottomrule
\end{tabular}
\end{table}

\begin{table}[H]
\centering
\caption{Per-location breakdown, isolated recordings only (\texttt{HS.csv}/\texttt{LS.csv}).}
\label{tab:perlocation}
\begin{tabular}{@{}llc c@{}}
\toprule
Component & Location & $n$ & Clipped \\
\midrule
Heart & LUSB & 13 & 0 \\
Heart & Apex & 10 & 0 \\
Heart & LLSB & 10 & 0 \\
Heart & RUSB & 7  & 0 \\
Heart & LC   & 6  & 0 \\
Heart & RC   & 4  & 0 \\
\addlinespace
Lung  & LUA & 11 & 0 \\
Lung  & RLA & 10 & 0 \\
Lung  & LMA & 9  & 0 \\
Lung  & LLA & 8  & 0 \\
Lung  & RUA & 7  & 0 \\
Lung  & RMA & 5  & 0 \\
\bottomrule
\end{tabular}
\end{table}

\textbf{Why this breakdown is usable.} Three properties make it trustworthy
as a basis for downstream decisions, rather than just a descriptive
aside. First, it is measured directly from the audio files themselves
(\texttt{wave} header + sample scan), not read off manifest metadata, so it
cannot inherit a CSV labeling error the way a metadata-only summary could.
Second, duration and sample rate being \emph{exactly} constant (std 0.0)
across every one of the 16 classes and 12 locations means no class or
location requires different preprocessing (resampling, padding, trimming)
before feature extraction --- a per-class confound that would otherwise have
to be ruled out separately for every downstream experiment. Third, the
clipped-file count is zero for \emph{every} class and \emph{every} location
here, which cross-checks the dataset-wide clipping finding
(Section~\ref{sec:audioquality}) at finer granularity: clipping is confined
entirely to \texttt{Mix.csv}'s \emph{mixed} component and does not
concentrate in any particular heart/lung class or recording site. That
matters because clipping correlated with a specific class would risk a
classifier learning the clipping artifact as a spurious class signal; here
that risk is ruled out by direct measurement rather than assumed.

\textbf{Why it is good news, with one caveat.} The absence of any per-class
or per-location quality defect means the planned classification stage
(\texttt{PROTOCOL.md} \S5.3) can treat every isolated recording as equally
trustworthy input --- there is no class whose recordings are disproportionately
degraded and would need to be excluded or reweighted on quality grounds
alone. The caveat is class \emph{size}, not quality: several heart classes
are very small (S4 $n{=}2$, AV Block and Tachycardia $n{=}3$ each), which
Table~\ref{tab:perclass} makes precise rather than leaving as an impression.
Under Section~\ref{sec:leakage}'s leak-group-respecting 5-fold split, a
class with only 2 recordings cannot appear in every fold's training set,
which will show up as unstable or undefined per-class accuracy for those
classes in the classification stage --- worth flagging now, before that
stage is built, rather than discovering it as a confusing result later
(Section~\ref{sec:conclusion}).

% ============================================================
\section{Recording Reuse and the Leakage Problem}
\label{sec:leakage}
% ============================================================

A na\"ive evaluation protocol would split \texttt{Mix.csv} by row (or by
\texttt{Mixed Sound ID}) into train/test folds, and fit any separation
dictionary on all of \texttt{HS.csv}/\texttt{LS.csv}. Both choices leak
ground truth into training, and the reason is specific to how HLS-CMDS was
assembled, not a generic caveat.

\subsection{Content-hash audit}

\texttt{HS.csv} and \texttt{Mix.csv} use \emph{different ID namespaces} for
the same underlying files (e.g.\ \texttt{HS.csv}'s \texttt{F\_N\_RC} vs.\
\texttt{Mix.csv}'s \texttt{H0001}), so file identity cannot be checked by ID
string comparison. Instead, \texttt{src/split.py} hashes the raw PCM samples
of every file (ignoring WAV headers, via 16-bit PCM read + MD5) to find:

\begin{itemize}
  \item \texttt{Mix.csv}'s 145 rows draw from only \textbf{89} distinct
    heart recordings and \textbf{74} distinct lung recordings (by content),
    each reused across multiple mix rows --- up to 6$\times$ for a single
    heart recording.
  \item \textbf{32} of those 89 heart recordings, and \textbf{30} of those
    74 lung recordings, are byte-identical (by PCM content) to a recording
    also listed standalone in \texttt{HS.csv}/\texttt{LS.csv} --- the exact
    pool a separation dictionary would be fit on.
\end{itemize}

\subsection{Leak groups}

Because recordings are reused \emph{across} mix rows independently for the
heart and lung components, the leakage-safe unit is not a single mix row but
a \emph{connected component}: two rows are linked whenever they share a
heart or a lung recording (directly or transitively through a third row).
Union-find over all 145 rows on this relation yields \textbf{29} connected
components (``leak groups''), ranging in size from 1 to 32 rows
(Table~\ref{tab:leakgroups}).

\begin{table}[H]
\centering
\caption{Leak-group size distribution over \texttt{Mix.csv}'s 145 rows (29 groups total; the sizes sum to 145). An earlier revision listed 13 groups of size one, which summed to 144 rows and 28 groups --- a transcription error, corrected against \texttt{split.py}'s output.}
\label{tab:leakgroups}
\begin{tabular}{@{}lcccccccccc@{}}
\toprule
Group size & 32 & 22 & 19 & 13 & 11 & 9 & 4 & 3 & 2 & 1 \\
\midrule
\# groups of this size & 1 & 1 & 1 & 1 & 1 & 1 & 2 & 3 & 4 & 14 \\
\bottomrule
\end{tabular}
\end{table}

\texttt{split.py}'s \texttt{assign\_folds()} partitions these 29 groups
(never splitting one across folds) into $k$ folds via a greedy
largest-group-to-smallest-fold assignment, giving reasonably balanced folds
despite the size-32 outlier (Table~\ref{tab:foldsizes}).
\texttt{dictionary\_pool()} then returns, for a given held-out fold, exactly
the \texttt{HS.csv}/\texttt{LS.csv} rows whose content does \emph{not}
appear anywhere in that fold's evaluation mixtures --- the pool any
dictionary-learning or model-fitting step must be restricted to when that
fold is held out.

\begin{table}[H]
\centering
\caption{5-fold split (seed 0): row counts and resulting dictionary-fitting pool sizes.}
\label{tab:foldsizes}
\begin{tabular}{@{}lccc@{}}
\toprule
Fold & Eval rows (of 145) & HS pool (of 50) & LS pool (of 50) \\
\midrule
0 & 32 & 44 & 46 \\
1 & 29 & 43 & 46 \\
2 & 28 & 44 & 47 \\
3 & 28 & 40 & 43 \\
4 & 28 & 47 & 38 \\
\bottomrule
\end{tabular}
\end{table}

This split governs \emph{both} dictionary/model fitting and evaluation ---
not evaluation alone --- which is the property \texttt{src/test\_split.py}
verifies directly (Section~\ref{sec:validation}).

% ============================================================
\section{Mixture Validity: the Additive-Triplet Audit}
\label{sec:validity}
% ============================================================

Section~\ref{sec:leakage}'s content-hash audit answers ``does a mix row's
ground truth leak into the pool used to fit a separation model.'' It says
nothing about a logically prior question: is the \emph{mixed} recording
named in a \texttt{Mix.csv} row actually related to the heart and lung
recordings named alongside it at all? \texttt{load\_dataset.py}'s
\texttt{verify\_mix\_triplets()} only checks that the three files exist, are
uniquely paired, and agree on format (Section~\ref{sec:dataset}) --- it
would pass even if \texttt{mixed} were a completely unrelated recording
bundled under the same row by a manifest error.

\texttt{verify\_additive\_triplets()} tests the assumption every separation
baseline and BSS Eval score in this pipeline actually relies on: $\text{mixed}
\approx a\cdot(\text{heart} + \text{lung})$ for some scalar gain $a$. For
each row it fits the least-squares $a$ minimizing
$\lVert \text{mixed} - a(\text{heart}+\text{lung})\rVert^2$, then reports the
relative residual $\lVert \text{mixed} - a\cdot\text{summed}\rVert /
\lVert \text{mixed}\rVert$. This measure is empirically bimodal with a wide
gap: genuinely additive rows land at residual $\sim 10^{-8}$--$10^{-4}$
(limited by 16-bit quantization), acoustically unrelated rows land at
residual $\sim 1$ (no shared energy at all), and no row falls in between ---
so the exact threshold (0.001) is not a sensitive tuning choice.

\textbf{Result: only 36 of 145 rows (24.8\%) pass.} The other 109 rows'
\texttt{mixed} file shares essentially no energy with its own named heart
and lung recordings. This independently reproduces
\emph{TEEP2026\_Sprint0\_Review}'s count on the same dataset copy, and (as
noted in Section~\ref{sec:audioquality}) the 36 valid rows are exactly the
36 rows showing the single-sample clipping fingerprint --- two independent
measurements (spectral/waveform residual, and a clipping artifact of the
mixing process) converging on the same row set is itself evidence the
finding is a real property of the data, not a bug in either check.

This has two consequences for the rest of this report. First, it is a
second, independent problem from leakage (Section~\ref{sec:leakage}): a row
can be leakage-safe and still be acoustically meaningless, or vice versa.
Second, because 75\% of \texttt{Mix.csv} is affected, every result in
Section~\ref{sec:results} is reported \emph{twice} --- on the full 145-row
set (for continuity with earlier work and Table~\ref{tab:foldsizes}'s
split) and on the 36-row valid subset alone (Table~\ref{tab:additivesubset})
--- since averaging rows that fail the additive-reference test into a headline separation-quality number
would mix ``the method failed to separate'' with ``there was nothing
coherent to separate in the first place.''

\begin{table}[H]
\centering
\caption{Additive-triplet audit result on HLS-CMDS as published
(Mendeley v3, provenance verified in Section~\ref{sec:provenance}).}
\label{tab:additivesubset}
\begin{tabular}{@{}lcc@{}}
\toprule
 & Rows & \% of 145 \\
\midrule
Additive (mixed $\approx a\cdot$(heart+lung)) & 36 & 24.8\% \\
Not additive (acoustically unrelated) & 109 & 75.2\% \\
\bottomrule
\end{tabular}
\end{table}

\textbf{Sensitivity of the verdict to the strictness of the model.} A
reviewer's natural objection is that a one-gain, zero-lag fit would reject a
mixture for trivial reasons --- a DC offset, an inverted polarity, a few
samples of misalignment, a different gain on each source. Each was tried on
all 145 rows (\texttt{additive\_audit\_sensitivity.py}): the mean removed
before fitting; the gain allowed to be negative (the least-squares gain is
already unconstrained, and 52 rows do choose a negative one --- none of
them passes at any threshold); the summed reference shifted by the
cross-correlation-optimal integer lag within $\pm100$ samples; separate
gains for heart and lung; and separate gains with separate per-source lags.
Under every one of the seven models the count at $10^{-3}$ is 36, with no
row gaining or losing the verdict. The passing cluster's residuals span
$0.9$--$4.5\times10^{-4}$ and the failing cluster's $0.94$--$1.0$ even under
the loosest model --- a gap of more than $2000\times$ --- so the thresholds
$10^{-2}$ and $10^{-3}$ give the same 36 rows, while $10^{-4}$ admits only
3, because 16-bit quantization alone puts a $\sim 10^{-4}$ floor on any
genuinely additive row. Among the 36, the optimal lag is 0 for every row;
among the 109 the optimal lag is spread uniformly to the $\pm100$-sample
edge, which is what cross-correlating two unrelated signals produces.
Every file reads 4000\,Hz. The 109 rows are therefore not near-misses of
the additive model under any relaxation tried; they are rows whose
mixture shares no energy with the two recordings named alongside it.

% ============================================================
\section{Methodology}
\label{sec:method}
% ============================================================

\subsection{BSS Eval metrics}
\label{sec:metrics}

Separation quality is scored with the classic BSS Eval decomposition
\cite{vincent2006}: an estimated source $\hat{s}_j$ is decomposed into a
target component (the best time-invariant-filtered version of the true
source $s_j$, filter length up to 512 taps), an interference component
(leakage from the other true sources), and an artifacts component (residual
distortion not explained by either). From this decomposition:

\begin{align}
\text{SDR} &= 10\log_{10}\frac{\lVert s_{\text{target}} \rVert^2}
  {\lVert e_{\text{interf}} + e_{\text{artif}} \rVert^2}
  & \text{(Signal-to-Distortion Ratio)} \\
\text{SIR} &= 10\log_{10}\frac{\lVert s_{\text{target}} \rVert^2}
  {\lVert e_{\text{interf}} \rVert^2}
  & \text{(Source-to-Interference Ratio)} \\
\text{SAR} &= 10\log_{10}\frac{\lVert s_{\text{target}} + e_{\text{interf}} \rVert^2}
  {\lVert e_{\text{artif}} \rVert^2}
  & \text{(Sources-to-Artifacts Ratio)}
\end{align}

\texttt{src/metrics.py} wraps \texttt{mir\_eval.separation.}
\texttt{bss\_eval\_sources} \cite{mireval}, specialized for this dataset's
2-source (heart, lung) case:

\begin{lstlisting}
def evaluate_heart_lung(heart_ref, lung_ref, heart_est, lung_est,
                         compute_permutation=True) -> dict:
    """Returns {"heart": {"sdr":.., "sir":.., "sar":..},
                "lung":  {"sdr":.., "sir":.., "sar":..}}"""
\end{lstlisting}

together with \texttt{evaluate\_mix\_row()}/\texttt{evaluate\_dataset()},
which load ground truth via \texttt{load\_dataset.py} and run a
caller-supplied \texttt{separate\_fn(mixed, sr)} $\to$
\texttt{(heart\_est, lung\_est)} over one or all \texttt{Mix.csv} rows, and
\texttt{summarize\_by\_class()} for per-class/per-location aggregation.
\texttt{mir\_eval} 0.8 deprecated \texttt{bss\_eval\_sources} in favor of an
unpublished museval-style API; since no working replacement exists yet, the
deprecated call is used with its \texttt{FutureWarning} suppressed locally.

\subsubsection{Validation}
\label{sec:validation}

\texttt{src/test\_metrics.py} (10 tests) validates the wrapper two ways:

\begin{itemize}
  \item \textbf{Cross-checks against \texttt{mir\_eval} directly} --- on
    synthetic arrays and on real dataset audio, the wrapper's output is
    compared bit-for-bit against a raw \texttt{mir\_eval} call on the same
    stacked arrays, catching wrapper bugs (stacking order, trimming,
    permutation labeling) independent of whether the underlying metric
    itself is correct.
  \item \textbf{Synthetic mixtures with closed-form ground truth.} White-noise
    sources are corrupted in ways whose correct BSS Eval score is
    analytically known, and the measured score must land within tolerance:

  \begin{table}[H]
  \centering
  \caption{Synthetic ground-truth validation (\texttt{test\_metrics.py}).}
  \label{tab:synthtests}
  \begin{tabular}{@{}llll@{}}
  \toprule
  Construction & Theoretical target & Tolerance & Result \\
  \midrule
  $\hat{s}=s$ (identity) & SDR/SIR/SAR $\to \infty$ & $>100$\,dB & $>300$\,dB \\
  $\hat{s}=s+n$, SNR 20\,dB & SDR $\approx$ SAR $\approx$ 20\,dB & $\pm 1.0$\,dB & within tolerance \\
  $\hat{s}=s+n$, SNR 6\,dB & SDR $\approx$ SAR $\approx$ 6\,dB & $\pm 1.0$\,dB & within tolerance \\
  $\hat{s}_j = s_j + \alpha s_k$, $\alpha{=}0.1$ & SIR $\approx -20\log_{10}\alpha = 20$\,dB & $\pm 1.0$\,dB & within tolerance \\
  $\hat{s}_j = s_j + \alpha s_k$, $\alpha{=}0.3$ & SIR $\approx 10.46$\,dB & $\pm 1.0$\,dB & within tolerance \\
  \bottomrule
  \end{tabular}
  \end{table}
\end{itemize}

\texttt{src/test\_split.py} (15 tests) validates the split, with the
load-bearing test being \texttt{test\_dictionary\_pool\_} \\
\texttt{excludes\_all\_held\_out\_recordings}: for every fold, it hashes
every recording in that fold's dictionary pool and asserts \emph{zero}
overlap with any recording used in that fold's held-out evaluation
mixtures. It also checks leak groups never split across folds, fold
coverage/balance, split reproducibility under a fixed seed, and that
\texttt{cross\_validate} (Section~\ref{sec:harness}) actually hands each
fold its restricted pool rather than the unrestricted full dataset. All 25
tests across both files pass ($\sim$2 minutes total, dominated by running
real BSS Eval over 15-second real recordings repeatedly).

\subsection{Evaluation harness}
\label{sec:harness}

\texttt{src/eval\_harness.py} ties Sections~\ref{sec:leakage}
and~\ref{sec:metrics} together:

\begin{lstlisting}
def cross_validate(fit_and_separate_fn, n_folds=5, seed=0):
    """For each fold: fit_and_separate_fn(hs_allowed, ls_allowed) -> separate_fn,
    evaluate separate_fn on the held-out fold via metrics.evaluate_dataset.
    Returns (results_df, fold_summary, cv_summary)."""
\end{lstlisting}

\texttt{fit\_and\_separate\_fn} is the caller's separation method or model;
the harness supplies each fold's leakage-restricted dictionary pool and
handles fold iteration, per-row scoring, and aggregation. An optional
\texttt{mix\_df} argument restricts evaluation to a caller-supplied subset
of \texttt{Mix.csv} (folds are then computed within that subset only) ---
used in Section~\ref{sec:results} to run every baseline on
Section~\ref{sec:validity}'s 36-row valid subset without duplicating the
fold-assignment logic. Results are aggregated three ways: mean SDR/SIR/SAR
per (fold, source); mean $\pm$ std \emph{across} folds per source
(\texttt{aggregate\_across\_folds}, fold-level dispersion); and mean,
median, and std pooled directly over every evaluated row per source
(\texttt{summarize\_pooled}, row-level dispersion). The last of these was
added after an observation in \emph{TEEP2026\_Sprint0\_Review}: averaging
within folds before taking a standard deviation can understate the true
row-to-row spread substantially (an across-fold figure of
$6.46\pm3.71$\,dB for the same data pools to $6.30\pm39.38$\,dB
row-by-row in one comparison the review made), so the two are not
interchangeable and a report should say which one it means. Both are
reported in Section~\ref{sec:results}.

\subsection{Synthetic mixture set and its own leakage-safe split}
\label{sec:synthetic}

Only 36 of \texttt{Mix.csv}'s 145 rows are additive
(Section~\ref{sec:validity}), too few to source a continuous
SDR-vs-difficulty sweep. \texttt{synthetic\_mix.py} builds a combinatorial
synthetic set instead: every one of the 50 heart recordings paired with
every one of the 50 lung recordings, mixed as
$\text{mixed} = a\cdot(\text{heart}+\text{lung}) + \text{noise}$ --- the same
functional form \texttt{verify\_additive\_triplets()}
(Section~\ref{sec:validity}) already shows the native additive rows follow,
not an arbitrary synthetic recipe. The gain $a$ is drawn log-uniformly per
source pair from the range empirically fitted on the 36 native additive
rows' own gains; additive Gaussian noise is scaled to a target SNR swept
over \texttt{SNR\_SWEEP\_DB} $= \{-5, 15, 35\}$\,dB (trimmed from 5 levels
after measuring that a 5-level sweep across all baselines would run
multi-hour) to give a controllable difficulty axis, rather than Han \&
Quan's own fixed 2\%-RMS noise level. \texttt{validate\_against\_native}
confirms the construction is faithful, not just claimed to be: rebuilding
each of the 36 native additive rows from their own fitted gain and zero
noise, then re-running \texttt{verify\_additive\_triplets()} against the
synthetic result, reproduces all 36 rows at the same $\sim 10^{-4}$ residual
the genuine native rows occupy.

\textbf{The split has to move from mixture row to source file.}
Section~\ref{sec:leakage}'s triplet-level split assigns \emph{leak
groups} --- connected components of mix rows sharing a heart or lung
recording --- to folds, which is the right unit when each recording appears
in only a few of the 145 native mixtures. Full combinatorial pairing breaks
that assumption: one heart file now appears in 50 mixtures, so within a few
hops every heart file transitively connects to every lung file through some
shared mixture, collapsing the entire synthetic set into a single connected
component with no split left to take. \texttt{assign\_source\_folds}
resolves this by assigning each of the 50 heart files and 50 lung files
\emph{independently} to one of 5 folds; a synthetic pair $(h,l)$ is only
used for held-out evaluation under fold $k$ when
$\text{fold}(h) = \text{fold}(l) = k$, and
\texttt{dictionary\_pool\_synthetic} excludes \emph{both} sources of every
held-out pair from that fold's dictionary-fitting pool --- the same
no-overlap contract as \texttt{split.dictionary\_pool}
(Section~\ref{sec:leakage}), re-derived at the file level rather than the
leak-group level. This is not a cosmetic change: getting it wrong (e.g.\
falling back to a row-level or leak-group split that no longer separates
anything) would let a fold's dictionary see a file that also appears in
that fold's held-out evaluation pairs, inflating every downstream SDR/SIR/SAR
number without any test failing to flag it. \texttt{test\_synthetic\_mix.py}
makes this a checked invariant rather than a one-off claim: for every fold,
\texttt{test\_dictionary\_pool\_excludes\_both\_sources\_of\_held\_out\_pairs}
asserts zero overlap between that fold's dictionary pool and either source
recording of any pair held out in that fold, alongside
\texttt{test\_only\_same\_fold\_pairs\_are\_included} confirming every
synthetic row actually satisfies $\text{fold}(h){=}\text{fold}(l)$.

\subsection{Baseline 1: bandpass filtering}
\label{sec:baseline1}

A fixed, non-adaptive separation method serves as the sanity floor any
learned method should comfortably beat --- if it doesn't, that signals a bug
in the learned method rather than a genuinely hard problem. Cutoffs were
derived from this dataset's own spectral content rather than assumed from
generic audio literature: Welch's method (\texttt{scipy.signal.welch},
1024-sample segments) was run over every isolated recording in
\texttt{HS.csv} and \texttt{LS.csv} separately and averaged per source,
showing heart energy dominant below $\sim$150--200\,Hz, lung energy dominant
from $\sim$200\,Hz up to $\sim$700--800\,Hz, and both sources falling into
the recording noise floor above $\sim$1\,kHz (Table~\ref{tab:psd} gives
representative values; no PSD plot is included in this draft).

\begin{table}[H]
\centering
\caption{Representative averaged, peak-normalized PSD values (dB) from the empirical audit.}
\label{tab:psd}
\begin{tabular}{@{}lccc@{}}
\toprule
Frequency & Heart (dB) & Lung (dB) & Dominant source \\
\midrule
93.8\,Hz  & $-2.3$  & $-12.1$ & Heart \\
187.5\,Hz & $-15.7$ & $-24.6$ & Heart (crossover region) \\
500\,Hz   & $-26.7$ & $-14.4$ & Lung \\
1000\,Hz  & $-50.4$ & $-50.2$ & Neither (noise floor) \\
\bottomrule
\end{tabular}
\end{table}

Based on this, Baseline~1 uses a heart band of 20--200\,Hz and a lung band
of 150--1000\,Hz, applied as a 4th-order zero-phase Butterworth bandpass
(\texttt{scipy.signal.butter} in second-order-section form, filtered with
\texttt{sosfiltfilt} to avoid the group-delay a single-pass IIR filter would
introduce). The 150--200\,Hz overlap between the two bands is intentional:
the PSD audit shows both sources carry genuine energy there, and forcing a
non-overlapping split would misrepresent the underlying physics as well as
erase the exact spectral-overlap failure mode this baseline exists to
expose. Filter order and type were \emph{not} tuned against the evaluation
data --- doing so would reproduce, in miniature, the leakage problem
Section~\ref{sec:leakage} exists to prevent, since it would fit the
``zero-training'' baseline to the very mixtures it is scored against.

\subsection{Baseline 2: supervised NMF}
\label{sec:baseline2}

Baseline~1 answers ``is the pipeline sane.'' It says nothing about whether
an actual published separation method works on this dataset, so
\texttt{src/baselines.py} also implements a leakage-safe reproduction of a
supervised non-negative matrix factorization (NMF) method from the
literature identified during related-work review
(\cite{nmfcrnn}; Section~\ref{sec:related}), following the classic
supervised-NMF separation recipe of Smaragdis (2007): two fixed
dictionaries, one per source, are learned once from isolated recordings via
KL-divergence multiplicative-update NMF (Lee \& Seung 2001), then frozen and
used to solve per-mixture activations against a held-out mixture.
Hyperparameters (dictionary ranks $K_i{=}10$ heart / $K_r{=}20$ lung, 100 MU
iterations to fit each dictionary, 60 MU iterations to solve activations)
were given during this project's method handoff as the paper's own values;
they have not yet been independently cross-checked against the primary
source (Section~\ref{sec:conclusion}). Separation is a soft, Wiener-style
mask applied to the mixture's complex STFT ($N_{\text{fft}}{=}512$,
hop 128), built from the two sources' NMF-reconstructed magnitude spectra.

\textbf{Leakage trap.} A na\"ive reproduction fits both dictionaries on
``all of \texttt{HS.csv}/\texttt{LS.csv}.'' Section~\ref{sec:leakage} shows
this lets a third of those recordings --- some of them a held-out mixture's
own ground truth --- into the dictionary the mixture is later scored
against. \texttt{make\_supervised\_nmf\_baseline()} instead receives only
each fold's leakage-safe \texttt{hs\_allowed}/\texttt{ls\_allowed} pool from
\texttt{eval\_harness.cross\_validate} (Section~\ref{sec:harness}) --- the
same restriction Baseline~1 is exposed to, though bandpass has no
dictionary to leak into. A leakage-safe score for this method is therefore
\emph{expected} to come in lower than a na\"ive reproduction's; that lower
number, not a higher one, is the methodologically correct result
(Section~\ref{sec:results}).

\subsection{Baseline 3: standard NMF (dictionary ablation)}
\label{sec:baseline3}

Baseline~2 conflates two effects: NMF factorization itself, and the
specific pretrained/frozen dictionary it factorizes against. Baseline~3
isolates the second effect by holding every other stage of Baseline~2's
pipeline fixed (same total rank $K_i{+}K_r{=}30$, same STFT parameters, same
Butterworth denoise pre-filter) while factorizing $W$ and $H$ directly out
of each held-out mixture's own spectrogram, with no dictionary-learning
phase from isolated recordings at all --- so, like Baseline~1, it never
touches \texttt{hs\_allowed}/\texttt{ls\_allowed} and has nothing to leak.
The resulting components are unlabeled by construction; they are assigned
to heart or lung post-hoc by spectral centroid against Baseline~1's own
PSD survey (heart energy concentrated below $\sim$200\,Hz), a fixed
physical prior rather than a source label, so the method stays genuinely
blind. The comparison this buys is direct: any gap between Baseline~2 and
Baseline~3 is attributable to the learned, frozen dictionary specifically,
not to NMF factorization or the shared preprocessing.

\subsection{Baseline 4: multi-stage Singular Spectrum Analysis (MSSA)}
\label{sec:baseline4}

Baseline~4 reproduces the two-stage MSSA method of Han \& Quan
\cite{ssa}, the same paper whose own Table~I is reported alongside this
project's results in Section~\ref{sec:results} for direct comparison. All
four hyperparameters are taken as stated in the paper rather than tuned:
window length $L{=}50$, cardiac/respiratory split at 250\,Hz (a
physiological frequency in absolute Hz, valid at this dataset's 4000\,Hz
sample rate since Nyquist, 2000\,Hz, clears it with margin), an eigenvalue
contribution threshold of $100/L = 2\%$, and a correlation threshold of
50\%. Stage 1 SSA-decomposes the raw mixture into 50 reconstructed
components (RC), routing each by its dominant Welch-PSD frequency into
either the cardiac output directly ($\le 250$\,Hz) or a residual pool
feeding Stage 2; Stage 2 re-decomposes that residual and selects
respiratory components by eigenvalue contribution, then Pearson correlation
against the already-selected set. Like Baseline~1, MSSA is zero-training ---
\texttt{hs\_allowed}/\texttt{ls\_allowed} go unused and no leakage trap
applies, since there is no fitting step at all.

\subsection{Baseline 5: EVMD (Enhanced Variational Mode Decomposition)}
\label{sec:baseline5}

Baseline~5 reproduces the EVMD-based lung-isolation stage of Puneet et al.\
\cite{edgelung}, the only related-work paper in Table~\ref{tab:related}
whose separation step reports no separation-quality metric at all despite
running directly on HLS-CMDS mixtures --- Table~\ref{tab:sixmethodsynth}
below is, to this project's knowledge, the first SDR/SIR/SAR computed for
that method on this dataset. Variational Mode Decomposition itself (Dragomiretzkiy
\& Zosso, 2014) is implemented from scratch (no maintained VMD package is
available), and the paper's own $K$-selection sweep ($K{=}2..10$ at
$\alpha{=}2000$, stopping once an energy-loss criterion and a per-mode
check both pass) is followed as stated. The four-criterion per-mode check
the paper describes (frequency-domain signature, NPE ratio, and a
Kurtosis-Index fallback) is under-specified in the primary source; this
project's best-effort reading of how the criteria combine is documented
inline in \texttt{code\_description.md} rather than presented as a literal
reproduction. Like Baselines~1 and~4, EVMD is zero-training. The full
$K$-sweep costs $\sim$15\,s/mixture, dominated by VMD's ADMM iterations;
Section~\ref{sec:results}'s synthetic-substrate column for this baseline
therefore runs on a stratified 75-row subsample rather than the full
1500-row synthetic set, a disclosed compute-driven scope reduction rather
than a silent one.

\subsection{Baseline 6: Conv-TasNet-lite}
\label{sec:baseline6}

Baselines~1--5 are either zero-training or a fixed-dictionary fit.
Baseline~6 (\texttt{convtasnet.py}) is this project's first neural,
end-to-end learned separation baseline, and the first place a training
loop with epochs and an optimizer appears in this codebase at all. It uses
Conv-TasNet's own encoder/TCN-separator/decoder design
\cite{luomesgarani2019}, sized down (``lite'', $\sim$325K parameters vs.\
Conv-TasNet's own $\sim$5M) for this dataset's much smaller training pool,
with a sigmoid mask per source (no unit-sum constraint, following
Conv-TasNet's own ablation finding that sigmoid performs at least as well as
a softmax constraint) and a fixed source order rather than
permutation-invariant training, since heart and lung are distinguishable
classes here rather than interchangeable speakers.

\textbf{Sample rate.} This dataset is natively 4000\,Hz; public
Conv-TasNet/Sepformer checkpoints are trained at 8--16\,kHz and do not
apply without resampling, and resampling up invents no real information.
The model is trained from scratch at the native 4000\,Hz instead --- not a
resolution limitation for this task, since heart (20--200\,Hz) and lung
(100--1000\,Hz) energy (Section~\ref{sec:baseline1}'s own PSD survey) sit
comfortably under 4000\,Hz's 2000\,Hz Nyquist.

\textbf{Training loop.} A fresh network is trained per fold from that
fold's leakage-safe \texttt{hs\_allowed}/\texttt{ls\_allowed} pool --- the
same contract every other baseline's \texttt{fit\_and\_separate\_fn}
receives (Section~\ref{sec:harness}). Since Conv-TasNet needs paired
(mixed, heart, lung) training examples rather than a spectral dictionary,
training mixtures are synthesized on the fly using
Section~\ref{sec:synthetic}'s own mixing recipe, so train-time mixtures are
drawn from the same distribution the model is evaluated against. A
file-level 80/20 split inside the allowed pool (never touching the outer
fold's held-out data) gives an internal validation set for early stopping
and learning-rate scheduling (AdamW, LR halved after 4 epochs without
validation-SI-SDR improvement, best-checkpoint restore).

\textbf{Honest scope.} This is the first neural baseline, establishing the
training-loop infrastructure --- not a tuned, converged model. The
dataset's $\sim$50 recordings/class (fewer per fold once leakage exclusions
apply) is small for a from-scratch neural separator even with on-the-fly
mixing augmentation. Rather than assume the shipped configuration
(\texttt{MAX\_EPOCHS=25}, \texttt{STEPS\_PER\_EPOCH=40}, \texttt{BATCH\_SIZE=8})
trained to convergence, \texttt{baseline6\_report.py} reports each fold's
actual training diagnostics and measured wall-clock directly
(Table~\ref{tab:b6diag}).

% ============================================================
\section{Results}
\label{sec:results}
% ============================================================

Per Section~\ref{sec:validity}, every baseline is reported on two row sets:
the full 145-row \texttt{Mix.csv} (Table~\ref{tab:baselinefull}, for
continuity with the fold sizes in Table~\ref{tab:foldsizes}) and the 36-row
additive-valid subset alone (Table~\ref{tab:baselinevalid}). All numbers are
5-fold cross-validated (\texttt{eval\_harness.cross\_validate}, seed 0),
reported as across-fold mean $\pm$ std (fold-level dispersion,
Section~\ref{sec:harness}); Table~\ref{tab:baselinevalid} additionally gives
the pooled row-level mean $\pm$ 95\% CI for the valid subset alongside the
fold mean, since the two conventions differ and both appear in this
report.

\subsection{Full 145-row set}

Table~\ref{tab:baselinefull} reports both baselines alongside a
no-separation reference (the raw mixed signal returned unmodified as both
the heart and lung estimate, i.e.\ what a classifier would see with no
separation step at all).

\begin{table}[H]
\centering
\caption{5-fold cross-validated separation quality on the full 145-row set: both baselines vs.\ no separation.}
\label{tab:baselinefull}
\begin{tabular}{@{}llccc@{}}
\toprule
Method & Source & SDR (dB) & SIR (dB) & SAR (dB) \\
\midrule
No separation & Heart & $-13.68 \pm 1.26$ & $1.71 \pm 0.93$  & $6.46 \pm 3.71$ \\
No separation & Lung  & $-15.89 \pm 0.84$ & $-1.55 \pm 0.91$ & $6.46 \pm 3.71$ \\
\addlinespace
Baseline 1 (bandpass) & Heart & $-12.65 \pm 1.40$ & $4.07 \pm 1.16$  & $-8.10 \pm 1.44$ \\
Baseline 1 (bandpass) & Lung  & $-13.83 \pm 1.12$ & $1.73 \pm 0.29$  & $-9.27 \pm 1.42$ \\
\addlinespace
Baseline 2 (supervised NMF) & Heart & $-12.82 \pm 1.31$ & $4.22 \pm 1.63$  & $-9.58 \pm 0.90$ \\
Baseline 2 (supervised NMF) & Lung  & $-15.09 \pm 0.98$ & $0.32 \pm 0.58$  & $-9.23 \pm 1.42$ \\
\bottomrule
\end{tabular}
\end{table}

\noindent (SAR is identical for heart and lung under ``no separation''
because with no processing applied, there is nothing to distinguish one
source's artifacts from the other's.)

Bandpass filtering improves SIR for both sources ($+2.4$\,dB heart,
$+3.3$\,dB lung) --- it does separate on frequency grounds --- but SAR
collapses by roughly 14--16\,dB relative to doing nothing, consistent with
filter ringing and destroyed harmonics dominating the residual error. SDR,
which folds interference and artifacts together, only nets a modest
$\sim$1--2\,dB gain over the no-separation reference as a result. This
matches the theoretical interference/artifact trade-off discussed in
Section~\ref{sec:relatedmetrics} and is exactly the outcome a bandpass
sanity-floor baseline is expected to produce: real but limited separation,
purchased at a real artifact cost. Baseline~2 (leakage-safe supervised NMF)
lands essentially tied with Baseline~1 on heart SDR (within one fold-std)
and $\sim$1.3\,dB worse on lung SDR --- a learned, source-specific
dictionary barely matching, and on lung losing to, a filter with zero
learned parameters.

\emph{But 109 of these 145 rows are non-additive under the named-reference test} (Section~\ref{sec:validity}):
for those rows there is no real heart+lung structure in \texttt{mixed} for
any method to recover, so a method's score on them measures how it responds
to essentially arbitrary audio, not separation quality. Table~\ref{tab:baselinefull}
is retained for continuity with the split reported in
Table~\ref{tab:foldsizes}, but the 36-row subset below is the number this
report treats as meaningful.

\subsection{36-row additive-valid subset}

Every number in Table~\ref{tab:baselinevalid} is computed from one file,
\texttt{results/canonical\_rows.csv} (\texttt{canonical\_rows.py}): the
BSS Eval scores of the cached separated waveforms of the S6-02 run
(\texttt{results/sdr\_sweep\_cache/}, 16-bit PCM, 2026-08-28), which are
also the exact audio Condition~B and every knee-point curve consume. Two
summaries of the same 36 rows are given side by side because both are in
use in this report and they differ: the across-fold mean of per-fold means
$\pm$ its std over the five folds (8/7/7/7/7 rows), and the pooled mean
over the 36 rows $\pm$ a 95\% $t$-interval over rows. An earlier revision
printed the fold-mean SDR here and the pooled SDR in Condition~B's table
($5.61$ vs.\ $5.46$\,dB for bandpass) without saying so; it also carried
a Baseline~2 entry ($3.78$/$-0.44$\,dB) from a run that predated the
primary-source corrections of Section~\ref{sec:baseline2}, and a
Baseline~6 entry ($5.15$\,dB) from a separate training run of the same
seed, which on the GPU is not bit-reproducible ($5.02$\,dB in the run
cached here). All three are superseded by this table.

\begin{table}[H]
\centering
\caption{Separation quality on the 36-row additive-valid subset, leakage-safe 5-fold, from \texttt{canonical\_rows.csv}. SDR is given under both conventions; SIR and SAR are pooled means. Baseline 0 scores the raw mixture as both estimates.}
\label{tab:baselinevalid}
\label{tab:pooledvalid}
\label{tab:sixmethodnative}
\small
\begin{tabular}{@{}llcccc@{}}
\toprule
 & & \multicolumn{2}{c}{SDR (dB)} & SIR & SAR \\
\cmidrule(lr){3-4}
Method & Source & fold mean $\pm$ std & pooled $\pm$ 95\% CI & (dB) & (dB) \\
\midrule
Baseline 0 (raw mixture) & Heart & $2.65 \pm 4.34$ & $2.48 \pm 2.79$ & $2.48$ & $74.6$ \\
Baseline 0 (raw mixture) & Lung & $-2.14 \pm 4.19$ & $-1.98 \pm 2.66$ & $-1.98$ & $74.6$ \\
\addlinespace
Baseline 1 (bandpass) & Heart & $5.61 \pm 3.68$ & $5.46 \pm 2.66$ & $6.30$ & $17.47$ \\
Baseline 1 (bandpass) & Lung & $3.94 \pm 3.57$ & $4.05 \pm 2.37$ & $5.73$ & $12.42$ \\
\addlinespace
Baseline 2 (supervised NMF) & Heart & $2.69 \pm 2.96$ & $2.58 \pm 1.87$ & $7.27$ & $7.58$ \\
Baseline 2 (supervised NMF) & Lung & $0.62 \pm 3.14$ & $0.71 \pm 2.34$ & $2.36$ & $9.93$ \\
\addlinespace
Baseline 3 (standard NMF) & Heart & $3.44 \pm 2.02$ & $3.35 \pm 1.88$ & $6.97$ & $9.37$ \\
Baseline 3 (standard NMF) & Lung & $2.22 \pm 3.91$ & $2.31 \pm 2.81$ & $5.02$ & $9.66$ \\
\addlinespace
Baseline 4 (MSSA) & Heart & $5.17 \pm 4.05$ & $5.02 \pm 2.91$ & $5.73$ & $18.81$ \\
Baseline 4 (MSSA) & Lung & $5.32 \pm 3.20$ & $5.39 \pm 1.87$ & $8.58$ & $12.18$ \\
\addlinespace
Baseline 5 (EVMD) & Heart & $3.73 \pm 3.43$ & $3.61 \pm 2.35$ & $4.99$ & $13.99$ \\
Baseline 5 (EVMD) & Lung & $0.36 \pm 3.39$ & $0.50 \pm 2.44$ & $1.40$ & $13.56$ \\
\addlinespace
Baseline 6 (Conv-TasNet-lite) & Heart & $5.19 \pm 4.14$ & $5.02 \pm 2.77$ & $6.28$ & $16.45$ \\
Baseline 6 (Conv-TasNet-lite) & Lung & $2.05 \pm 3.20$ & $2.13 \pm 2.21$ & $4.77$ & $9.07$ \\
\bottomrule
\end{tabular}
\end{table}

Two things stand out on the valid subset. \textbf{First, absolute separation
quality is far higher than Table~\ref{tab:baselinefull} suggests}: SDR jumps
by 15--18\,dB for both methods and both sources once rows with no genuine
mixture structure are excluded, confirming that the full-145-row numbers
were dominated by rows where nothing coherent could be separated, not by a
uniformly hard separation problem. \textbf{Second, and more consequentially,
the ranking between the two baselines reverses relative to
Table~\ref{tab:baselinefull}}: on the valid subset, zero-training bandpass
\emph{beats} leakage-safe supervised NMF by $\sim$2.9\,dB SDR on heart and
$\sim$3.3\,dB SDR on lung (fold means), a materially larger and more one-sided gap than
the near-tie observed on the full (mostly non-additive) row set. A learned,
source-specific separation method failing to beat a fixed filter with zero
learned parameters, on exactly the rows where genuine separation is
possible, is not the result a na\"ive reading of the reproduced paper would
predict, and is discussed further in Section~\ref{sec:discussion}.

\subsection{Six methods, three substrates}
\label{sec:sixmethod}

Sections~\ref{sec:baseline3}--\ref{sec:baseline6} add four more baselines,
and Section~\ref{sec:synthetic} adds the 1500-row synthetic set as this
project's primary evaluation substrate going forward, with the 36-row
native-additive subset retained alongside as a distinct column rather than
folded in --- agreement or disagreement between the two substrates is
itself audit evidence, not noise to average away. All numbers below are
mean $\pm$ 95\% CI; Baseline~5's synthetic column runs on a stratified
75-row subsample (Section~\ref{sec:baseline5}), not the full 1500.

\begin{table}[H]
\centering
\caption{Synthetic substrate (primary), 5-fold, $n{=}1500$ except Baseline~5 ($n{=}75$).}
\label{tab:sixmethodsynth}
\begin{tabular}{@{}llccc@{}}
\toprule
Method & Source & SDR (dB) & SIR (dB) & SAR (dB) \\
\midrule
Baseline 1 (bandpass) & Heart & $2.10 \pm 0.43$ & $6.17 \pm 0.47$ & $11.01 \pm 0.45$ \\
Baseline 1 (bandpass) & Lung  & $-1.24 \pm 0.44$ & $4.25 \pm 0.40$ & $5.92 \pm 0.48$ \\
\addlinespace
Baseline 2 (supervised NMF) & Heart & $-0.70 \pm 0.40$ & $6.54 \pm 0.44$ & $3.47 \pm 0.34$ \\
Baseline 2 (supervised NMF) & Lung  & $-3.73 \pm 0.40$ & $1.39 \pm 0.39$ & $3.59 \pm 0.38$ \\
\addlinespace
Baseline 3 (standard NMF) & Heart & $-0.00 \pm 0.43$ & $6.93 \pm 0.50$ & $4.84 \pm 0.37$ \\
Baseline 3 (standard NMF) & Lung  & $-3.53 \pm 0.46$ & $4.06 \pm 0.39$ & $1.39 \pm 0.46$ \\
\addlinespace
Baseline 4 (MSSA) & Heart & $1.59 \pm 0.47$ & $5.58 \pm 0.47$ & $11.48 \pm 0.52$ \\
Baseline 4 (MSSA) & Lung  & $-1.60 \pm 0.50$ & $6.24 \pm 0.42$ & $3.74 \pm 0.55$ \\
\addlinespace
Baseline 5 (EVMD) & Heart & $2.42 \pm 1.70$ & $6.20 \pm 2.02$ & $10.15 \pm 1.48$ \\
Baseline 5 (EVMD) & Lung  & $-2.99 \pm 2.25$ & $2.81 \pm 1.90$ & $5.63 \pm 2.53$ \\
\addlinespace
Baseline 6 (Conv-TasNet-lite) & Heart & $3.16 \pm 0.40$ & $6.47 \pm 0.51$ & $12.14 \pm 0.31$ \\
Baseline 6 (Conv-TasNet-lite) & Lung  & $0.37 \pm 0.35$ & $4.29 \pm 0.41$ & $6.96 \pm 0.27$ \\
\bottomrule
\end{tabular}
\end{table}

The native-additive column of this comparison is Table~\ref{tab:baselinevalid}
above (its pooled convention matches the synthetic table's).

\textbf{Substrate agreement.} Every baseline scores lower on the synthetic
substrate than on the 36-row native set (heart gaps $-1.2$ to $-3.4$\,dB,
lung gaps $-1.8$ to $-7.0$\,dB, computed as synthetic pooled mean minus native pooled
mean), consistently in the same direction across all six methods. The
synthetic set's much larger $n$ makes its own CI tight regardless of this
gap, so a tight synthetic-column CI is not itself evidence the two
substrates agree --- the gap column is the evidence, and it says the
synthetic set is harder, not that either substrate is unreliable. The
lung-side gap is consistently larger than the heart-side gap for every
method, which is consistent with Baseline~1's own PSD survey
(Section~\ref{sec:baseline1}): lung energy is more spread and closer to the
noise floor, so it is more sensitive to the synthetic set's noise sweep
than heart energy is.

\textbf{Baseline 6 still does not clearly win.} Conv-TasNet-lite scores
highest of the six methods on synthetic heart and lung SDR, and highest on
native-additive heart SDR, but on native-additive lung SDR it trails
Baseline~4 (MSSA, $5.39$ vs.\ $2.14$\,dB) by more than either method's CI
half-width. A learned, end-to-end model that only sometimes beats a
zero-training method is the expected outcome of an intentionally
undertrained first neural baseline (Section~\ref{sec:baseline6}), not
evidence the architecture is unsuited to the task; Table~\ref{tab:b6diag}
below is what determines whether ``undertrained'' is actually the right
explanation.

\begin{table}[H]
\centering
\caption{Baseline 6 training diagnostics, synthetic-set folds (\texttt{max\_epochs=25}, \texttt{steps\_per\_epoch=40}, \texttt{batch\_size=8}).}
\label{tab:b6diag}
\begin{tabular}{@{}ccccc@{}}
\toprule
Fold & Epochs trained & Best val.\ SI-SDR (dB) & Device & Fit time (s) \\
\midrule
0 & 25 & $0.78$ & cuda & 27.1 \\
1 & 25 & $1.78$ & cuda & 26.2 \\
2 & 25 & $2.65$ & cuda & 26.2 \\
3 & 25 & $3.35$ & cuda & 26.3 \\
4 & 25 & $1.68$ & cuda & 26.2 \\
\bottomrule
\end{tabular}
\end{table}

Every fold trains for the full 25-epoch budget without the LR-halving
criterion (4 epochs without validation-SI-SDR improvement) ever firing ---
best-validation-SI-SDR is still improving, or at worst plateauing briefly,
when training stops, which reads as compute-budget-limited rather than
converged. Combined with only $\sim$32 training recordings per class per
fold (Table~\ref{tab:foldsizes}-scale pool sizes, further split 80/20 for
internal validation), this supports Section~\ref{sec:baseline6}'s own
``honest scope'' framing: Baseline~6 establishes that the training-loop
infrastructure works end-to-end (Table~\ref{tab:sixmethodsynth} shows finite,
non-degenerate output), not that this architecture has been given a fair
chance to demonstrate its ceiling on this dataset.

% ============================================================
\section{Downstream Classification: Does Separation Hurt Accuracy?}
\label{sec:classifier}
% ============================================================

Every section so far measures \emph{separation quality} (SDR/SIR/SAR).
None of it yet answers the question the whole project exists to answer
(Section~\ref{sec:yaqub}): does separating heart sound from a mixture,
before classifying it, actually cost accuracy, and if so, how much, at
what separation quality? Answering that needs a classifier, a way to
compare its accuracy on clean vs.\ separated audio under a controlled
split, and --- because separation quality is continuous, not a single
number --- a way to sweep separation quality and trace out a curve rather
than compare two fixed points.

\subsection{Why one classifier architecture, deliberately, at first}
\label{sec:onearchitecture}

The classifier in this section is described in the project charter as a
\emph{measuring instrument}, not a research contribution in its own
right: its only job is to reveal whether separation quality and accuracy
are related, not to be the best possible heart-sound classifier. That
framing is the reason exactly one architecture was built first
(Section~\ref{sec:robustness} revisits this once the classifier's own
result turns out to matter more than originally planned) --- a second
architecture only becomes necessary once a reader might reasonably ask
``is this result about separation, or about the one classifier you
happened to pick,'' which is a question about the \emph{measuring
instrument}, not about classification research for its own sake.

\subsection{The class-grouping decision}
\label{sec:classgrouping}

HLS-CMDS labels each isolated heart recording with one of ten specific
types (Table~\ref{tab:perclass}). Classifying all ten directly is not
possible under this project's leak-group split: S4 has only 2 recordings
total, and under a 5-fold split built from leak groups (Section~\ref{sec:leakage})
there is no guarantee those 2 recordings land in different folds --- if
they land in the same fold, three of the five folds train with \emph{zero}
S4 examples at all, and AV Block and Tachycardia ($n{=}3$ each) are barely
better off. This is exactly the small-class risk Section~\ref{sec:perclass}
flagged before this section was built.

Two fixes were considered: collapse to a binary Normal/Abnormal
classification, or group the ten types into a small number of clinically
meaningful categories. The categories chosen --- Normal, Murmur, Extra
Sound, Rhythm Disorder --- are not an arbitrary compromise: they are the
\emph{exact} four-class scheme Yaqub et al.\ (Section~\ref{sec:yaqub})
already use for this same dataset, confirmed by reading their primary
source directly (their Tables~6--9), and reproduce the same class list
their own 89\%$\to$41\% collapse is measured on. Mapping HLS-CMDS's ten
types onto those four is unambiguous from standard cardiology terminology:
Normal stays Normal ($n{=}9$); the four murmur sub-types (Mid/Late
Systolic, Late Diastolic, Early Systolic) become Murmur ($n{=}24$); S3 and
S4 --- both extra heart sounds occurring outside the normal S1/S2 pair ---
become Extra Sound ($n{=}7$); and Atrial Fibrillation, Tachycardia, and
AV Block --- all disturbances of the heart's electrical rhythm or
conduction, as opposed to its valves --- become Rhythm Disorder
($n{=}10$). This solves the small-class problem (the smallest group,
Extra Sound at $n{=}7$, is comfortably large enough to appear in every
fold) and, unlike collapsing to binary Normal/Abnormal, keeps this
project's accuracy numbers directly comparable to the 86--89\% figures
Yaqub et al.\ report on the very scheme this project's motivating result
is anchored to.

\subsection{Architecture 1: MFCC summary statistics + an RBF-kernel SVM}
\label{sec:arch1}

Given the class-grouping decision above and Section~\ref{sec:onearchitecture}'s
``measuring instrument, not a contribution'' framing, the classifier itself
was chosen to be small, well-understood, and unlikely to overfit a dataset
this size (50 recordings total, as few as $\sim$35--40 in a fold's training
split): 13 Mel-Frequency Cepstral Coefficients (MFCCs, a standard,
perceptually-motivated compression of a spectrogram into a small number of
coefficients per time frame, widely used in speech and audio
classification before deep learning), pooled to a single 26-dimensional
feature vector per recording by taking the mean and standard deviation of
each coefficient over time, fed into a Support Vector Machine (SVM) with a
Radial Basis Function (RBF) kernel and class-balanced loss weighting (so
the majority Murmur class does not simply dominate the optimization by
volume alone). An SVM with a fixed, moderate number of support vectors is
far less prone to memorizing a training set this small than a
gradient-trained deep model would be, which is precisely the trade-off
Section~\ref{sec:robustness} tests directly with a second architecture
rather than assuming.

Evaluation follows the same leak-group folds every separation baseline
uses (\texttt{assign\_hs\_folds()}, extending \texttt{split.py}'s existing
fold assignment from mixture rows to the isolated \texttt{HS.csv}
recordings themselves), stratified by class group for any recording never
reused inside a mixture, so no fold is starved of a minority class. Table~
\ref{tab:conditionAfold} reports the resulting 5-fold cross-validated
result --- the first accuracy number in this report; every prior result is
a dB-scale SDR/SIR/SAR figure.

\begin{table}[H]
\centering
\caption{Condition A (Architecture 1, MFCC+SVM): isolated-heart-sound classification on the 50 \texttt{HS.csv} recordings, 5-fold cross-validated. Not the same evaluation set as Condition~B's ``isolated'' row (Table~\ref{tab:conditionb}), which scores the same per-fold classifiers on the 36 native-additive mixtures' own heart references; the two accuracies ($60.0\%$ vs.\ $52.8\%$) differ by design.}
\label{tab:conditionAfold}
\begin{tabular}{@{}lcc@{}}
\toprule
Metric & Mean $\pm$ 95\% CI ($n{=}5$ folds) & \\
\midrule
Accuracy & $60.0\% \pm 10.7\%$ & \\
Macro-F1 & $0.43 \pm 0.12$ & \\
\bottomrule
\end{tabular}
\end{table}

With only 5 folds, this confidence interval should be read as indicative,
not decisive (the same caveat Section~\ref{sec:harness} already applies to
every separation baseline's across-fold dispersion) --- it says roughly
where the accuracy sits, not that a repeat run would land within
$\pm 0.1$ percentage points of it.

\subsection{Failure-mode analysis: the confusion matrix}
\label{sec:classifier-failure}

An aggregate accuracy number hides \emph{which} classes are actually being
confused for which. Section~\ref{sec:bg-classmetrics}'s confusion matrix
makes this concrete: pooling every fold's held-out predictions (each
recording predicted exactly once, so pooling never double-counts), the
per-class recall is Normal 44\%, Murmur 92\%, Extra Sound 29\%, Rhythm
Disorder 20\%. The dominant confusion is directional, not symmetric: the
minority classes' errors pull toward Murmur --- Extra Sound is predicted
Murmur \emph{more} often than it is predicted correctly (43\% vs.\ 29\%
recall), and Rhythm Disorder is predicted Murmur or Extra Sound (30\% each)
more often than itself (20\%) --- while Murmur is almost never mistaken
for anything else (0\% confused as Normal, 4\% each as Extra Sound and
Rhythm Disorder). Normal is the one exception, whose errors run toward
Rhythm Disorder (33\%) instead. This is consistent with Murmur's
training-set dominance ($n{=}24$ of 50) pulling the SVM's decision
boundary toward it despite class-balanced loss weighting --- reweighting
the training loss does not guarantee balanced \emph{predictions} once the
feature space itself gives the four classes limited separability to begin
with.

\subsection{A controlled way to sweep separation quality: the degradation scheme}
\label{sec:degradation}

Comparing accuracy at exactly two points --- perfectly isolated audio, and
one separation method's real output --- is what Yaqub et al.\ do
(Section~\ref{sec:yaqub}), and the charter's own C2 research question asks
for more: not just \emph{whether} separation costs accuracy, but the
\emph{knee point} --- the separation quality (SDR) below which separation
starts actively hurting accuracy relative to not separating at all. A knee
point cannot be identified from two endpoints; it needs a continuous curve
across a range of separation quality.

\texttt{synthetic\_mix.py}'s noise sweep (Section~\ref{sec:synthetic})
already varies difficulty \emph{before} separation runs, by controlling
how noisy the input mixture is. A second, complementary knob varies
difficulty \emph{after} separation: take a real separation method's actual
output $s$ and its ground truth $g$, and blend them,
\begin{equation}
\text{degraded}(\alpha) = (1-\alpha)\,g + \alpha\, s = g + \alpha\,(s-g),
\qquad \alpha\in[0,1],
\label{eq:degrade}
\end{equation}
so $\alpha{=}0$ reproduces the clean ground truth exactly and $\alpha{=}1$
reproduces that method's real separated output exactly. This has the same
``signal plus a scaled error term'' shape as the noise sweep in
Equation~\ref{eq:linearmix}'s synthetic-mixing recipe, except the injected
term $s-g$ is a \emph{real separation method's own characteristic error}
(its own specific mix of interference and artifact, Section~\ref{sec:bg-bsseval})
rather than generic Gaussian noise --- the two knobs are kept side by side
deliberately, because they answer different questions: the noise sweep
asks ``how does accuracy degrade as the separation \emph{problem} gets
harder,'' while the interpolation sweep asks ``at a given separation
quality, however it was reached, how does accuracy behave,'' holding one
specific method's error \emph{shape} fixed and varying only its
magnitude.

\textbf{A real, measured caveat about this scheme, found rather than
assumed.} Validating Equation~\ref{eq:degrade} on real audio revealed that
a \emph{uniform} grid of $\alpha$ values does \emph{not} give an even
spread of measured SDR: on the row checked, SDR fell from
$\sim$262\,dB at $\alpha{=}0$ (near-bit-identical to the clean signal) to
$\sim$10.6\,dB by $\alpha{=}0.1$ alone, then flattened out to
$\sim$4.8\,dB by $\alpha{=}1$ --- because decibels are logarithmic and most
sensitive to a small \emph{absolute} error right where the signal is
nearly clean. A uniform $\alpha$ grid therefore wastes most of its points
on the flat tail. The first implementation therefore solved for $\alpha$
by bisection, on the assumption --- checked on one row, not proved --- that
SDR is monotone in $\alpha$. That assumption was later audited on every
curve and found to fail for seven of them, and the bisection has been
replaced; both are described here because the sweep's validity rests on
it.

\textbf{SDR($\alpha$) has a closed form, so the audit is exhaustive.}
Equation~\ref{eq:degrade} is linear in $\alpha$, and BSS Eval's
decomposition of an estimate into target and error components is a fixed
linear projection onto the span of the delayed references. Every
component of $d(\alpha)$ is therefore the same convex combination of the
components of $g$ and of $s$, and
$\text{SDR}(\alpha) = 10\log_{10} N(\alpha)/D(\alpha)$ with $N$ and $D$
quadratics in $\alpha$ whose six coefficients are inner products of the two
decompositions. Two BSS Eval decompositions per (method, row, source) give
the exact curve at any resolution (\texttt{sdr\_alpha\_audit.py}); the
closed form was checked against \texttt{mir\_eval} at five interior
$\alpha$ values on 72 curves and agrees to $10^{-5}$\,dB wherever
\texttt{mir\_eval}'s permutation search keeps the heart/lung assignment.
On a 1001-point grid, 425 of the 432 curves (six methods, 36 rows, two
sources) are monotone non-increasing to within $10^{-6}$\,dB. The seven
that are not are all Conv-TasNet-lite heart curves: the network's
estimate on those rows is close to a polarity-inverted copy of the
reference, so the blend passes near cancellation and SDR dips to between
$-6$ and $-20$\,dB at $\alpha \approx 0.04$--$0.17$ before rising to the
method's real output ($-7$ to $+19$\,dB). For those rows targets below the
real output have two roots, and the bisection, which only ever probed
$\alpha \ge 10^{-3}$ on a log grid and assumed one crossing, clamped them
to the real output --- a silently wrong point. On the monotone curves the
bisection was adequate: median $|\alpha_{\text{bisection}} -
\alpha_{\text{exact}}| = 2\times10^{-4}$, 90th percentile
$3.5\times10^{-3}$, and the achieved SDR within $0.17$\,dB of the target
at the 90th percentile.

\textbf{The sweep now solves for $\alpha$ exactly, under a pre-stated
rule.} Since $\text{SDR}(\alpha) = t$ is the quadratic
$N(\alpha) - 10^{t/10} D(\alpha) = 0$, every root in $[0,1]$ is found in
closed form, with no grid resolution to worry about even where the curve
falls $230$\,dB in the first $10^{-3}$ of $\alpha$. The rule, fixed before
the re-run: take the \emph{smallest} root --- the crossing on the branch
that starts at the ground truth --- and record the number of roots and
whether the row's curve is monotone in the provenance. A target with no
root in $[0,1]$ (above SDR(0), or below the curve's minimum, which for a
monotone curve is the method's own output) is recorded as unattainable and
the point is placed at $\alpha{=}1$ with its clamp flagged, as before; it
is never searched for. Each sweep point's SDR is scored with the
heart/lung assignment fixed, because inside a cancellation dip the heart
estimate carries almost no energy and \texttt{mir\_eval}'s permutation
search would score the lung estimate as ``heart'' instead; the permuted
value is stored alongside and flagged where it differs. Re-generating the
provenance from the cached separated audio under this rule changed the
$\alpha$ of 99 of 1053 attainable heart targets by more than $0.01$ --- 22 of
them on the seven non-monotone rows, the rest where the bisection's
$0.1$\,dB tolerance or a permutation flip in its measured SDR had moved it
--- and un-clamped 55 points; every attainable point now sits at its
target to four decimals;
every curve, knee and bootstrap figure in this report is from the
re-generated sweep, and the point grid remains
$\{25, 20, 15, 10, 5, 0, -5\}$\,dB.

\subsection{Generating the SDR sweep}
\label{sec:sweepgen}

For every one of the six separation baselines, every one of the 36 native
additive rows (Section~\ref{sec:validity}), and both sources, the sweep
generator (\texttt{sdr\_sweep.py}) fits that baseline once per fold (the
same leakage-safe contract every earlier baseline uses,
Section~\ref{sec:harness}), separates the held-out mixture, caches the
real separated audio to disk, and then applies Section~\ref{sec:degradation}'s
target-SDR root-finding to hit every value in the target grid --- one
separation call per row, then seven cheap degraded points reconstructed
from it, rather than seven expensive re-separations. This produces up to
$6 \times 36 \times 2 \times 7 = 3024$ provenance rows, each recording
which baseline, which mixture, which source, what SDR was targeted, what
$\alpha$ actually achieved it, and the resulting measured SDR/SIR/SAR.

Because this involves re-running every separation baseline --- including
the two most expensive, EVMD (Section~\ref{sec:baseline5}, $\sim$15\,s per
mixture) and Conv-TasNet-lite (Section~\ref{sec:baseline6}, which needs a
freshly trained network per fold) --- generation is scripted to checkpoint
incrementally and resume safely if interrupted, rather than needing to
run start-to-finish in one sitting. The generation run is complete: all
$3024$ provenance rows exist (six baselines $\times$ 36 rows $\times$ two
sources $\times$ seven targets, five folds each), and the cache-then-
reconstruct path was re-checked on 20 sampled rows against the recorded
achieved SDR with a maximum discrepancy of $0.25$\,dB. One property of the
sweep matters for reading every curve below: the target grid runs from
$-5$ to $25$\,dB, but a point cannot be degraded \emph{below} the
baseline's own real output, so targets under that floor are clamped to it
(29--37\% of rows per baseline). Each baseline's curve therefore starts at
its own real SDR --- $2.6$\,dB for Baseline~2, $5.5$\,dB for Baseline~1
--- and its two lowest target points coincide with its Condition~B
operating point.

\subsection{Condition B: classifying separated audio, and the weight-sharing decision}
\label{sec:conditionb}

A subtlety has to be settled before ``classify the separated audio'' is
well-defined: is the classifier used on separated audio \emph{the same
trained weights} as Condition A, or a fresh classifier retrained on
separated training audio? These are different experiments answering
different questions --- the former isolates ``does separation hurt a
classifier that never saw separation artifacts during training,'' the
latter asks ``can a classifier adapt to separation artifacts if it is
allowed to.'' This project uses the \emph{same trained weights} per fold,
matching Yaqub et al.'s own methodology exactly (their Experiment~3 model
\emph{is} their Experiment~4 model, with no retraining in between,
Section~\ref{sec:yaqub}): reproducing their finding under controlled
conditions means matching their causal claim, and same-weights is also
the cleaner isolation of the separation method as the only variable
between the two conditions.

A second, easy-to-miss consistency requirement follows: the classifier's
fold assignment and the sweep's fold assignment must agree on what
``fold~$k$'' means, or a fold's classifier could unknowingly be evaluated
on separated audio derived from a recording it was partly trained on.
Because \texttt{sdr\_sweep.py} computes its folds over the 36-row
native-additive subset specifically (not the full 145-row \texttt{Mix.csv}
Condition~A's own standalone fold assignment defaults to), the classifier's
fold assignment is explicitly rebuilt from that same 36-row basis for
Condition~B, rather than assumed to already match.

\textbf{A confound found in the first run, and removed.} The first
complete Condition~B run produced a result that looked like a finding and
was a defect: Baselines~1, 4 and~6 --- a fixed filter, an SVD
decomposition and a neural network --- returned the \emph{same 36
predictions}. The cached separated audio was checked first and was
distinct per method (different content hashes, different RMS), and the
loading path resolves each baseline to its own directory, so the
classifier was genuinely seeing different audio and returning the same
answer. The reason was level. \texttt{HS.csv}'s recordings span
$0.0005$--$0.013$ RMS, while every baseline's heart estimate inherits the
mixture's peak-normalising gain and sits at $0.03$--$0.09$ RMS ---
$20$--$25\times$ louder than anything the classifier had trained on.
MFCC$_0$ is log-energy, so separated inputs landed at a standardised
distance of $11$--$15$ from the origin against a training cloud whose
farthest point is at $7.4$; an RBF kernel evaluates to $\approx 0$ that far
from every support vector, and the SVM's decision function collapsed to
the same vector ($[0.95, 3.13, -0.16, 2.10]$ for bandpass and Conv-TasNet
alike) for every far-away input. The prediction was a per-fold constant
--- Rhythm Disorder for fold~0's classifier, Murmur for the other four
--- and did not depend on the audio. Worse, the sweep's blend
$(1-\alpha)\,\text{ref} + \alpha\,\text{est}$ moved loudness back toward
the training range as SDR rose, so the first accuracy-vs-SDR curve was
partly a loudness-vs-SDR curve.

The fix is to make both classifiers invariant to absolute level: every
waveform is scaled to unit RMS before feature extraction
(\texttt{heart\_classifier.level\_normalize}, applied identically by
both backends at training and at inference). Nothing else changed. After
it, isolated references sit inside the training cloud (mean standardised
magnitude $\approx 1$ per feature) and separated estimates remain outside
it --- at distance $8$--$13$, now on the spectral-shape coefficients
(MFCC$_1$--MFCC$_4$) rather than on energy --- which is the distortion
Condition~B exists to measure: a 20--200\,Hz bandpass removes most of the
spectral envelope a heart-sound classifier trained on full-band
recordings relies on. Condition~A also changed with the retrained
classifiers (Table~\ref{tab:conditionAfold} reports the post-fix values);
every Condition~B and knee-point number in this report is from the
corrected run, and the pre-fix run is kept only in this paragraph.

Table~\ref{tab:conditionb} reports the result for all six baselines on
the same 36 rows, with the same per-fold weights throughout.

\begin{table}[H]
\centering
\caption{Condition B, both architectures, 36 native additive rows, same per-fold weights as Condition A. The isolated row is those 36 rows' own heart references scored with the same per-fold classifiers --- a different, smaller set than Condition~A's 50 recordings (Table~\ref{tab:conditionAfold}), which is why it reads $52.8\%$ here and $60.0$/$62.0\%$ there. Accuracy is pooled over the 36 rows; every interval is a 95\% percentile interval from 5000 cluster-bootstrap replicates over the 24 leak groups the rows fall into (\texttt{condition\_b\_stats.py}). $\Delta$ is isolated minus separated accuracy, paired by row; $\Delta_{\text{raw}}$ is separated minus raw-mixture accuracy, the knee question at one point. Mean SDR is each method's real heart output (Table~\ref{tab:baselinevalid}, pooled).}
\label{tab:conditionb}
\small
\begin{tabular}{@{}llcccc@{}}
\toprule
Arch.\ & Input & Accuracy [95\% CI] & SDR & $\Delta$ vs.\ isolated [95\% CI] & $\Delta_{\text{raw}}$ [95\% CI] \\
\midrule
1 & Isolated ground truth & $52.8$ [$35.3$, $70.0$] & --- & --- & --- \\
1 & Raw mixture (no separation) & $38.9$ [$24.1$, $53.1$] & $2.48$ & $+13.9$ [$0.0$, $29.0$] & --- \\
1 & Baseline 1 (bandpass) & $30.6$ [$14.3$, $45.2$] & $5.46$ & $+22.2$ [$2.5$, $45.2$] & $-8.3$ [$-28.6$, $8.1$] \\
1 & Baseline 2 (supervised NMF) & $41.7$ [$24.3$, $57.9$] & $2.58$ & $+11.1$ [$-8.6$, $33.3$] & $+2.8$ [$-17.1$, $20.0$] \\
1 & Baseline 3 (standard NMF) & $27.8$ [$14.3$, $39.5$] & $3.35$ & $+25.0$ [$11.6$, $41.9$] & $-11.1$ [$-24.1$, $-2.4$] \\
1 & Baseline 4 (MSSA) & $36.1$ [$19.4$, $51.2$] & $5.02$ & $+16.7$ [$-2.4$, $38.9$] & $-2.8$ [$-20.6$, $11.4$] \\
1 & Baseline 5 (EVMD) & $27.8$ [$10.7$, $43.2$] & $3.61$ & $+25.0$ [$3.2$, $48.6$] & $-11.1$ [$-32.4$, $7.4$] \\
1 & Baseline 6 (Conv-TasNet-lite) & $36.1$ [$21.2$, $51.4$] & $5.02$ & $+16.7$ [$-2.6$, $38.2$] & $-2.8$ [$-23.5$, $16.2$] \\
\addlinespace
2 & Isolated ground truth & $52.8$ [$34.3$, $69.2$] & --- & --- & --- \\
2 & Raw mixture (no separation) & $33.3$ [$17.9$, $50.0$] & $2.48$ & $+19.4$ [$-2.7$, $38.9$] & --- \\
2 & Baseline 1 (bandpass) & $16.7$ [$3.1$, $30.6$] & $5.46$ & $+36.1$ [$19.4$, $52.8$] & $-16.7$ [$-34.5$, $-2.6$] \\
2 & Baseline 2 (supervised NMF) & $44.4$ [$29.0$, $57.6$] & $2.58$ & $+8.3$ [$-10.5$, $25.0$] & $+11.1$ [$-6.9$, $27.8$] \\
2 & Baseline 3 (standard NMF) & $27.8$ [$11.8$, $44.4$] & $3.35$ & $+25.0$ [$7.4$, $40.5$] & $-5.6$ [$-22.9$, $12.1$] \\
2 & Baseline 4 (MSSA) & $16.7$ [$3.1$, $30.6$] & $5.02$ & $+36.1$ [$19.4$, $52.8$] & $-16.7$ [$-34.5$, $-2.6$] \\
2 & Baseline 5 (EVMD) & $16.7$ [$3.1$, $30.6$] & $3.61$ & $+36.1$ [$19.4$, $52.8$] & $-16.7$ [$-34.5$, $-2.6$] \\
2 & Baseline 6 (Conv-TasNet-lite) & $36.1$ [$19.4$, $52.9$] & $5.02$ & $+16.7$ [$-3.4$, $34.4$] & $+2.8$ [$-18.4$, $23.5$] \\
\bottomrule
\end{tabular}
\end{table}

The direction matches Yaqub et al.'s collapse for every one of the six
methods under both classifiers, and the intervals say how much of it is
resolvable at $n{=}36$. Under Architecture~1 the drop is $11$--$25$ points,
and the leak-group bootstrap interval of the paired difference excludes
zero for three methods (bandpass, standard NMF, EVMD); a group-level
sign-flip permutation test gives $p = 0.075$, $0.004$ and $0.062$ for the
three, and $0.18$--$0.43$ for the rest. Under Architecture~2 the drop is
$8$--$36$ points and excludes zero for four (bandpass, standard NMF, MSSA,
EVMD; permutation $p = 0.003$--$0.031$). The earlier fold-level paired
$t$-tests ($n{=}5$) are retained in \texttt{condition\_b\_report.html}
but are not the basis of any claim: five fold means of unequal size are a
weaker and less honest unit than 36 paired rows resampled by leak group.
The $\Delta_{\text{raw}}$ column asks the knee question at one point ---
is the separated audio classified \emph{better than the untouched
mixture}? --- and the answer is no for every method under both
classifiers: every interval includes zero or lies below it, and for
standard NMF (Architecture~1) and for bandpass, MSSA and EVMD
(Architecture~2) the whole interval is below zero.

Two things in the table run against a naive reading of the SDR column.
The \emph{smallest} drop belongs to the method with the \emph{lowest}
SDR: supervised NMF, at $2.6$\,dB, loses $11$ (Architecture~1) and $8$
points (Architecture~2), where bandpass, at $5.5$\,dB, loses $22$ and
$36$. And the three methods that strip spectral content --- bandpass, MSSA
and EVMD --- produce nearly the same predictions as one another under the
SVM ($31$--$34$ of 36 identical) and \emph{exactly} the same 36
predictions under the CNN, which scores all three at $16.7\%$. That last
number is a per-fold constant again: the CNN's logits on those inputs
differ by method but share the same arg-max in every fold (Extra Sound in
folds 0 and 3, Normal in the others), so a classifier trained on
full-band recordings, faced with audio from which most of the spectral
envelope has been removed, falls back to a default class whatever the
method. Level normalisation removed one cause of this collapse
(Section~\ref{sec:conditionb}); the spectral one is not a defect but the
phenomenon Condition~B measures, and it means that at their real output
the three stripping methods are indistinguishable to either classifier
--- the accuracy differences between them in the table are differences of
zero to three rows. The classifier is evidently responding to how much
of the spectral envelope survives separation, which SDR, a
waveform-domain energy ratio, does not measure directly. This table
reproduces the direction of Yaqub's result at roughly half its size; the
trajectories below are what the C2 question actually needs.

\subsection{The accuracy-vs-SDR curve and the knee point}
\label{sec:knee}

Once every target-SDR point in the sweep (Section~\ref{sec:sweepgen}) has
been classified with Condition~B's weight-sharing setup, the result is a
curve per baseline: measured SDR on the horizontal axis, classification
accuracy on the vertical axis. This is, as far as this project has found,
the first time HLS-CMDS classification accuracy and separation-quality SDR
have been placed on the same axis at all --- Yaqub et al.\ report accuracy
with no SDR attached to either endpoint, and every separation-quality
table in this report (Section~\ref{sec:results}) reports SDR with no
downstream accuracy attached.

\textbf{Definition of the knee point.} Rather than an arbitrary
curvature-based ``elbow'' heuristic, the knee point is defined directly
from the charter's own wording: the SDR at which a baseline's accuracy
curve crosses the accuracy of \emph{not separating at all} (classifying
the raw, unseparated mixture directly with the same classifier weights),
walking from high SDR down to low SDR. Above the knee, separating helps
(or at least does not hurt); below it, a classifier would have done as
well or better without separating at all. The crossing point is found by
linear interpolation between the two adjacent \emph{measured} points that
bracket it --- not a smoothed curve fit through all seven points, since
smoothing through only seven points (each an average over a handful of
rows) would manufacture a precision the data does not support. Four
outcomes are possible and are reported distinctly rather than collapsed
into one number: a clean crossing; the curve staying above the
no-separation reference across the whole measured range (no knee found,
separation always helps here); staying below across the whole range
(separation never helps); or a ``noisy'' crossing in the wrong direction,
flagged as such rather than silently treated as a clean knee.

\textbf{The curve.} Figure~\ref{fig:knee} is the complete six-baseline
curve for Architecture~1, all $1512$ heart-source sweep points classified
with Condition~B's per-fold weights; Table~\ref{tab:knee} gives the knee
points the detector found on it.

\begin{figure}[H]
\centering
\includegraphics[width=\textwidth]{../results/plots/sdr_knee_point_curve_heart_classifier.png}
\caption{Classification accuracy against measured SDR, Architecture~1
(MFCC+SVM), 36 native additive rows per point. The dashed line is the
same classifier on the raw mixture ($38.9\%$); the dash-dot line is
isolated ground truth ($52.8\%$). Dotted verticals mark each baseline's
knee point. Each curve's leftmost points sit at that baseline's own real
output SDR, because the sweep cannot degrade below it.}
\label{fig:knee}
\end{figure}

\begin{table}[H]
\centering
\caption{Raw-mixture crossings, Architecture 1: the SDR at which a baseline's trajectory, walking from high to low SDR, crosses the no-separation reference ($38.9\%$), by linear interpolation between the two bracketing sampled points; ``always above'' means no sampled point falls below the reference. Bootstrap columns are from Table~\ref{tab:kneeboot} (row resampling); ``n.i.'' = not identifiable within the 25\,dB grid. Real output SDR is the method's own output (Table~\ref{tab:baselinevalid}, pooled).}
\label{tab:knee}
\begin{tabular}{@{}lcccc@{}}
\toprule
Baseline & Crossing (dB) & $P$(crossed) & 95\% interval & Real output SDR (dB) \\
\midrule
Baseline 1 (bandpass)         & $20.1$ & $0.71$ & [$7.3$, n.i.] & $5.46$ \\
Baseline 2 (supervised NMF)   & always above & $0.43$ & --- & $2.58$ \\
Baseline 3 (standard NMF)     & $15.0$ & $0.90$ & [$5.0$, n.i.] & $3.35$ \\
Baseline 4 (MSSA)             & $15.6$ & $0.72$ & [$7.3$, n.i.] & $5.02$ \\
Baseline 5 (EVMD)             & $18.2$ & $0.85$ & [$6.1$, n.i.] & $3.61$ \\
Baseline 6 (Conv-TasNet-lite) & $4.5$  & $0.63$ & [$3.1$, n.i.] & $5.02$ \\
\bottomrule
\end{tabular}
\end{table}

\textbf{What the trajectories show.} The six split into two groups that
the SDR axis alone would not predict. Bandpass, standard NMF, MSSA and
EVMD (Baselines~1, 3, 4, 5) all start at $25$--$33\%$ at their own real
output --- below the $38.9\%$ no-separation reference --- and climb to
$44$--$53\%$ at $25$\,dB, crossing the reference between $15$ and
$20$\,dB. Supervised NMF and Conv-TasNet-lite (Baselines~2 and~6) start
\emph{at or above} the reference despite having the lowest and a
middling SDR respectively: Baseline~2 never falls below it at any sampled
point and sits at the isolated ceiling from $10$\,dB upward, and
Baseline~6's only excursion below it is its two lowest points ($36.1\%$
against $38.9\%$, one row), which the detector reports as a crossing at
$4.5$\,dB. Both are methods whose output is the mixture's own STFT under
a magnitude mask (or a learned resynthesis of it), so the spectral
envelope the classifier was trained on survives. Bandpass, MSSA and EVMD
discard spectral content by construction; standard NMF is a mask method
too, but its mask is fitted on the mixture alone with no source
dictionary, and under this classifier it behaves like the three that
strip. For those four the figure says what it said before the level fix,
now without the confound: at the operating point each actually achieves,
its output classifies $6$--$14$ points worse than the untouched mixture,
and along this interpolation protocol it would need roughly $15$ to
$20$\,dB of heart-source SDR --- $10$ to $15$\,dB more than it delivers
--- before separating beat not separating. For the two mask-based methods
the same classifier is indifferent to separation at every SDR: their
output never hurts and, at the low SDR they reach, never helps either.

\textbf{What they do not show.} The vertical resolution is one row in 36
($2.8$ points), and a binomial 95\% interval on a 36-row accuracy near
$40\%$ is about $\pm16$ points, wider than the gap between adjacent
trajectories at most SDRs. The bootstrap columns of Table~\ref{tab:knee}
make this concrete: for every baseline the upper end of the crossing's
interval is the top of the measured grid, and for bandpass, MSSA and
Conv-TasNet a quarter to a third of resamples have no crossing at all.
Individual crossing values within the late-crossing group are therefore
not distinguishable from one another, and Baseline~6's $4.5$\,dB is a
one-row excursion that the CNN below does not reproduce. What survives
the noise is the grouping of Baselines~2 and~6 against Baselines~1, 4
and~5 --- consistent across all seven target levels and, as
Section~\ref{sec:robustness} shows, across both classifiers --- and the
fact that no spectral-stripping method's operating point is anywhere
near its crossing region. The crossing \emph{region} for that group,
roughly $15$--$20$\,dB under this classifier, is the quotable figure; the
individual values are not.

\subsection{Robustness check: does a second, different classifier see the
same knee?}
\label{sec:robustness}

Section~\ref{sec:onearchitecture} explained why one classifier
architecture was acceptable while separation quality (not accuracy) was
this project's headline result. That reasoning no longer applies once the
knee point itself \emph{is} the headline result: the first question a
skeptical reader will ask is whether a knee point is a property of
separation quality, or an artifact of the one architecture that happened
to measure it. The only way to answer that is to repeat the entire curve
with a second, architecturally distinct classifier and see whether the two
knee points agree.

``Architecturally distinct'' matters here, not just ``a second model'':
two different SVMs, or an SVM with different hyperparameters, would not
isolate the right variable, since both would still share the same feature
representation and the same kind of decision boundary. The second
architecture built for this check instead uses a log-Mel spectrogram (a
2-D time-frequency image, the same kind of representation Baseline~6's
encoder learns implicitly, Section~\ref{sec:bg-convtasnet}) fed into a
small convolutional neural network (two convolutional blocks, global
average pooling, a single linear output layer) --- a gradient-trained
model operating on a fundamentally different input shape than Architecture~1's
pooled MFCC statistics, chosen specifically to maximize the chance that a
disagreement, if found, reflects something real about the two
architectures' different decision boundaries rather than two flavors of
the same idea.

\begin{table}[H]
\centering
\caption{Architecture 2 (log-Mel + CNN) vs.\ Architecture 1 (MFCC+SVM), Condition A, 5-fold cross-validated.}
\label{tab:arch2}
\begin{tabular}{@{}lcc@{}}
\toprule
Metric & Architecture 1 (MFCC+SVM) & Architecture 2 (log-Mel+CNN) \\
\midrule
Accuracy & $60.0\% \pm 10.7\%$ & $62.0\% \pm 3.9\%$ \\
Macro-F1 & $0.43 \pm 0.12$    & $0.54 \pm 0.11$ \\
\bottomrule
\end{tabular}
\end{table}

\textbf{The first configuration was underfitting, and the fix was
optimisation, not architecture.} Architecture~2's first configuration ---
raw log-Mel input, 40 full-batch Adam steps --- scored $30.0\%$, far below
Architecture~1, with a confusion matrix collapsed toward \emph{Normal}.
An earlier revision read that as the small-data overfitting risk of
Section~\ref{sec:arch1} coming true. It was the opposite: the model's
\emph{training} accuracy was $23$--$40\%$ per fold, so it had not fitted
the training data either. Two causes, both mundane. The log-Mel features
sit around $-12.4 \pm 2.3$ (a $10^{-6}$ floor on quiet recordings), an
offset the first convolution had to absorb before learning anything; and
40 full-batch steps is 40 gradient updates in total. Standardising each
Mel bin with the training fold's own mean and standard deviation (stored
on the classifier, applied identically at test time, so nothing from a
held-out fold enters), training in mini-batches of eight for 300 epochs
with weight decay $10^{-3}$, and changing nothing about the network,
brings training accuracy to $85$--$95\%$. Three configurations were run
in total and all three are recorded in the module's docstring; the
adopted one was chosen as the first whose training accuracy cleared
$85\%$ on every fold --- the point at which underfitting was demonstrably
gone --- and the reader should know it also had the highest test accuracy
of the three at the time, and weigh that accordingly. That tuning was
done before the level normalisation of Section~\ref{sec:conditionb} was
added to both backends; with it, and no further change, Architecture~2
scores $62.0\% \pm 3.9\%$ (fold accuracies $60$--$70\%$, the tightest
spread of any classifier in this report) with macro-F1 $0.54$. Its
confusion matrix no longer collapses: Normal recall $67\%$, Murmur
$71\%$, Extra Sound $71\%$, Rhythm Disorder $30\%$ --- a more balanced
profile than Architecture~1's, which buys its $92\%$ Murmur recall with
$20$--$29\%$ on the minority classes.

\textbf{The second curve.} Figure~\ref{fig:knee2} repeats the entire
sweep classification with the tuned Architecture~2, on the same 36-row
fold basis, and Table~\ref{tab:kneecompare} sets the two architectures'
knee points side by side.

\begin{figure}[H]
\centering
\includegraphics[width=\textwidth]{../results/plots/sdr_knee_point_curve_heart_classifier_cnn.png}
\caption{The same sweep classified with Architecture~2 (log-Mel + CNN).
Its no-separation reference is $33.3\%$ and its isolated accuracy on
these 36 rows is $52.8\%$.}
\label{fig:knee2}
\end{figure}

\begin{table}[H]
\centering
\caption{Knee points under both architectures. ``always above'' means the curve never falls below that architecture's no-separation reference at any measured SDR, so no knee exists in range. The detector's agreement criterion is both knees clean and within $3$\,dB.}
\label{tab:kneecompare}
\begin{tabular}{@{}lccc@{}}
\toprule
Baseline & Arch.\ 1 knee (dB) & Arch.\ 2 knee (dB) & Spread (dB) \\
\midrule
Baseline 1 (bandpass)         & $20.1$ & $9.3$ & $10.7$ \\
Baseline 2 (supervised NMF)   & always above & always above & --- \\
Baseline 3 (standard NMF)     & $15.0$ & always above & --- \\
Baseline 4 (MSSA)             & $15.6$ & $9.4$ & $6.3$ \\
Baseline 5 (EVMD)             & $18.2$ & $7.9$ & $10.3$ \\
Baseline 6 (Conv-TasNet-lite) & $4.5$  & always above & --- \\
\bottomrule
\end{tabular}
\end{table}

\textbf{Verdict: the knee \emph{value} is not robust to the architecture;
the grouping and the shape are.} Under the study's 3-dB operational
agreement criterion --- a working threshold chosen for this analysis, not
a standard --- no baseline agrees: the three crossings both architectures
find are $6$--$11$\,dB apart, every one lower under the CNN, and for the
other three the CNN finds no crossing where the SVM finds one or, for
supervised NMF, neither does. Two things account for most of that. The
reference bar is lower for Architecture~2 ($33.3\%$ against $38.9\%$),
so the same trajectory crosses it earlier; and the CNN's trajectories are
steeper and cleaner --- every one of its six is monotone or nearly so
across the seven levels, rising from $17$--$42\%$ to $53$--$61\%$, where
the SVM's are flat for the first three levels and only then climb. A
crossing quoted to the decibel is therefore a property of the classifier
that measured it, on this dataset at this sample size, and this report
does not quote one.

What the two sets of trajectories share is the part that matters for the
C2 question, and it is more than shape. Both split the six methods the
same way: supervised NMF is ``always above'' the reference under both,
Conv-TasNet is always above under the CNN and dips below by one row under
the SVM, standard NMF is always above under the CNN and crosses at
$15.0$\,dB under the SVM, and the three spectral-stripping methods ---
bandpass, MSSA, EVMD --- cross late under both, at $15$--$20$\,dB for the
SVM and $8$--$9$\,dB for the CNN. Under both, those three methods'
operating points ($3.6$--$5.5$\,dB) lie below their crossing region, so
their output classifies worse than the raw mixture. The robustness check
therefore weakens the numerical claim of Section~\ref{sec:knee} and
sharpens its qualitative one: what a downstream classifier responds to is
not SDR as such but whether the separator preserves the spectral
envelope, and the methods that do not would need to improve by several
decibels before separating beat not separating.

\textbf{How uncertain each crossing is, measured.} Table~\ref{tab:kneeboot}
bootstraps the whole construction --- trajectory, SDR axis and reference
together, 2000 replicates, once over rows and once over the 24 leak
groups the 36 rows belong to --- and finds the crossing on every replicate
with the same detector (\texttt{sdr\_knee\_bootstrap.py}). Two things
the point estimates could not show. First, a crossing is not even
guaranteed to exist: for Conv-TasNet under Architecture~1 a third of
replicates find none, and for the mask methods under Architecture~2 most
do not. Second, where it does exist its interval spans most of the
measured range: every Architecture~1 interval reaches the $25$\,dB top of
the grid and is reported as not identifiable there, and the tightest
interval in the table (EVMD under the CNN, $3$--$15$\,dB) is still
$12$\,dB wide. The cluster bootstrap moves nothing by more than about a
decibel, so the rows' dependence is not what limits the precision;
$n{=}36$ is. What is stable across replicates is the ordering the text
above relies on: under the CNN, Baselines~1, 4 and~5 cross in
$90$--$98\%$ of replicates and Baseline~2 in $15\%$, so ``the strippers
have a crossing and supervised NMF does not'' survives resampling even
though no individual crossing value does.

\begin{table}[H]
\centering
\caption{Bootstrap uncertainty of the raw-mixture crossings, 2000 replicates, resampling the 36 rows (i.i.d.) or the 24 leak groups they fall into (cluster). Trajectory, its SDR axis and the no-separation reference are all recomputed on each replicate and the crossing found with the same detector. $P$(crossed) is the fraction of replicates with a clean crossing; the interval is the 2.5--97.5 percentile over those replicates; ``n.i.'' (not identifiable) means the percentile sits at the $25$\,dB top of the measured grid.}
\label{tab:kneeboot}
\small
\begin{tabular}{@{}llcccc@{}}
\toprule
 & & \multicolumn{2}{c}{Rows ($n{=}36$)} & \multicolumn{2}{c}{Leak groups ($n{=}24$)} \\
\cmidrule(lr){3-4}\cmidrule(lr){5-6}
Arch.\ & Baseline & $P$(crossed) & Median [95\% CI] & $P$(crossed) & Median [95\% CI] \\
\midrule
1 & B1 bandpass        & $0.71$ & $15.4$ [$7.3$, n.i.] & $0.75$ & $15.6$ [$8.5$, n.i.] \\
1 & B2 supervised NMF  & $0.43$ & $6.2$ [$3.4$, 22.5] & $0.39$ & $5.7$ [$3.1$, 18.7] \\
1 & B3 standard NMF    & $0.90$ & $15.0$ [$5.0$, n.i.] & $0.97$ & $15.0$ [$6.0$, 23.8] \\
1 & B4 MSSA            & $0.72$ & $15.4$ [$7.3$, n.i.] & $0.75$ & $15.6$ [$7.3$, n.i.] \\
1 & B5 EVMD            & $0.85$ & $18.0$ [$6.1$, n.i.] & $0.92$ & $18.1$ [$7.3$, n.i.] \\
1 & B6 Conv-TasNet     & $0.63$ & $18.4$ [$3.1$, n.i.] & $0.65$ & $19.0$ [$2.8$, n.i.] \\
\addlinespace
2 & B1 bandpass        & $0.95$ & $9.5$ [$4.3$, 21.7] & $0.98$ & $9.5$ [$3.6$, 22.5] \\
2 & B2 supervised NMF  & $0.15$ & $3.9$ [$1.7$, 22.5] & $0.15$ & $3.8$ [$1.7$, n.i.] \\
2 & B3 standard NMF    & $0.45$ & $4.2$ [$1.7$, 13.7] & $0.45$ & $4.6$ [$1.7$, 14.2] \\
2 & B4 MSSA            & $0.94$ & $9.4$ [$4.3$, 21.7] & $0.94$ & $9.3$ [$3.5$, 23.4] \\
2 & B5 EVMD            & $0.90$ & $8.3$ [$3.2$, 15.2] & $0.93$ & $8.2$ [$2.5$, 15.2] \\
2 & B6 Conv-TasNet     & $0.27$ & $8.2$ [$2.9$, 22.5] & $0.25$ & $8.5$ [$2.0$, 23.3] \\
\bottomrule
\end{tabular}
\end{table}


% ============================================================
\section{Computational Cost}
\label{sec:compute}
% ============================================================

Every result so far asks ``how good'' a method is. A separate, equally
practical question --- especially relevant given this project's own
related-work review of an FPGA-deployed lung classifier (\cite{edgelung},
Table~\ref{tab:related}) --- is ``how expensive'' each method is to
actually run. This section measures that directly, using
Section~\ref{sec:bg-compute}'s vocabulary (MACs, latency, the Pareto
front), for all six separation baselines plus the Condition~A classifier.

\subsection{MACs per method}
\label{sec:macs}

Table~\ref{tab:macs} reports the number of multiply-accumulate operations
each method performs to process one real 60{,}000-sample (15\,s) mixture,
alongside each method's persisted parameter count (zero for every
zero-training method). Two precision levels are reported explicitly rather
than presented as equally certain: \emph{exact} counts come directly from
a method's own literal matrix-multiplication or convolution-layer shapes
(or, for Conv-TasNet-lite, from forward hooks recording real output shapes
during an actual forward pass, rather than hand-deriving stride/padding/dilation
arithmetic by hand); \emph{estimate} counts (MSSA's SVD, EVMD's ADMM loop)
use standard textbook complexity formulas applied to this project's real
problem sizes, and are order-of-magnitude figures, not literal
op-for-op tallies.

\begin{table}[H]
\centering
\caption{MACs and parameters per method, one real 60,000-sample mixture/recording.}
\label{tab:macs}
\begin{tabular}{@{}llll@{}}
\toprule
Method & Parameters & MACs/inference & Precision \\
\midrule
Baseline 1 (bandpass) & 0 & 2.4\,M & exact \\
Condition A classifier (MFCC+SVM) & 40 support vectors & 8.1\,M & exact \\
Baseline 2 (supervised NMF) & 7{,}710 (frozen dictionary) & 219\,M & exact \\
Baseline 3 (standard NMF) & 0 (transductive) & 1.16\,G & exact \\
Baseline 4 (MSSA) & 0 & 1.20\,G & estimate \\
Baseline 6 (Conv-TasNet-lite) & 325{,}465 & 1.96\,G & exact \\
Baseline 5 (EVMD) & 0 & 5.22\,G & estimate \\
\bottomrule
\end{tabular}
\end{table}

Two findings are worth calling out specifically. First, Conv-TasNet-lite's
measured parameter count (325{,}465) independently reproduces the
$\sim$325K figure Section~\ref{sec:baseline6} already reports from
PyTorch's own model summary --- a genuine cross-check between two
independently-computed numbers, not a coincidence of rounding. Second,
Baseline~3 costs $\sim$5.3$\times$ more MACs per mixture than Baseline~2
despite being the ``simpler'' ablation (no learned dictionary,
Section~\ref{sec:baseline3}): Baseline~3 is fully \emph{transductive}
(every one of its 160 multiplicative-update iterations updates both $W$
\emph{and} $H$ fresh per mixture), while Baseline~2 only solves
activations (60 iterations, $H$ only) against an already-frozen
dictionary. This gives the inductive-vs-transductive distinction already
noted in Section~\ref{sec:baseline3} a concrete computational-cost number
to go with the qualitative one.

\subsection{Desktop latency under a uniform protocol}
\label{sec:latency}

MACs are a theoretical proxy (Section~\ref{sec:bg-compute}); the actual
number a deployment decision cares about is wall-clock time. Latency was
measured under a \emph{uniform protocol}, held identical across every
method regardless of what else is running on the machine: one untimed
warm-up call (so any one-time cache or JIT warm-up cost is excluded), then
several timed repetitions on the same real 60{,}000-sample input, with the
\emph{median} (not mean, which a single slow repetition could dominate)
and interquartile range reported for both wall-clock time and process CPU
time. One-time setup cost --- NMF dictionary fitting, Conv-TasNet-lite
training --- is fit once, untimed, before the timed region begins, since
this section measures \emph{inference} latency, matching the convention
Yaqub et al.\ themselves use for their own model comparison (their
Table~5 reports inference time, not training time).

\textbf{Two runs, and why the first is retracted.} The first measurement
(2026-08-27) was taken on a shared machine under a 1-minute load average
of 74.6 on 16 cores --- this project's own sweep generation plus other
users' jobs --- and was reported with that caveat and the argument that a
uniform protocol keeps the \emph{ranking} meaningful. The re-measurement
below, on the same machine idle (load average $0.6$ at start, $1.6$ at
end), shows that argument was wrong: the loaded run inflated the
CPU-bound methods by between $10\times$ (bandpass, $36 \to 3.6$\,ms) and
$340\times$ (standard NMF, $21{,}036 \to 61$\,ms) while leaving the
GPU-bound Conv-TasNet-lite nearly untouched ($23.9 \to 16.7$\,ms), so it
did not preserve the ranking either --- it reordered the two fastest
methods. Every latency-derived claim in this report is from the idle run.
Table~\ref{tab:latency} gives it with the dispersion a reader needs to
judge it; the measurement environment is recorded by the script itself
(Table~\ref{tab:latencyenv}).

\begin{table}[H]
\centering
\caption{Desktop inference latency per 15\,s clip (60{,}000 samples), uniform protocol: one untimed warm-up, then $n$ timed repetitions; median, IQR and p5--p95 of wall-clock time, median process CPU time, and the real-time factor ($15\,\text{s}/\text{median}$; above 1 is faster than real time). Per-call preprocessing inside each method (e.g.\ Baseline~2's denoising bandpass and STFT) is included; one-time fitting and training are not. Machine idle: load average $0.6$/16 cores.}
\label{tab:latency}
\small
\begin{tabular}{@{}lccccc@{}}
\toprule
Method & $n$ & Wall median (ms) & IQR / p5--p95 (ms) & CPU median (ms) & Real-time factor \\
\midrule
Baseline 1 (bandpass)            & 5 & $3.63$    & $0.01$ / $3.62$--$3.70$     & $3.6$    & $4{,}130\times$ \\
Condition A classifier (MFCC+SVM) & 5 & $3.69$   & $0.02$ / $3.67$--$3.74$     & $63.6$   & $4{,}060\times$ \\
Baseline 6 (Conv-TasNet-lite), GPU & 5 & $16.7$  & $0.03$ / $16.7$--$16.7$     & $16.7$   & $899\times$ \\
Baseline 2 (supervised NMF)      & 5 & $17.2$    & $2.2$ / $17.1$--$22.2$      & $107$    & $873\times$ \\
Baseline 3 (standard NMF)        & 5 & $61.4$    & $0.7$ / $60.6$--$61.6$      & $367$    & $244\times$ \\
Baseline 4 (MSSA)                & 5 & $1{,}097$ & $39$ / $1{,}057$--$1{,}117$ & $7{,}519$ & $13.7\times$ \\
Baseline 5 (EVMD)                & 3 & $16{,}812$ & $49$ / $16{,}794$--$16{,}882$ & $16{,}812$ & $0.89\times$ \\
\bottomrule
\end{tabular}
\end{table}

\begin{table}[H]
\centering
\caption{Measurement environment for Table~\ref{tab:latency}, recorded at run time (\texttt{results/latency\_environment.csv}).}
\label{tab:latencyenv}
\small
\begin{tabular}{@{}ll@{}}
\toprule
CPU & Intel Core i9-9900K @ 3.60\,GHz, 16 logical CPUs \\
GPU & NVIDIA GeForce RTX 3060 (Conv-TasNet-lite only; all other methods CPU) \\
RAM & 65.7\,GB \\
OS / Python & Linux 6.8.0 (glibc 2.39) / 3.12.4 \\
Libraries & numpy 2.4.6, scipy 1.17.1, librosa 0.11.0, scikit-learn 1.9.0, torch 2.3.1+cu121, mir\_eval 0.8.2 \\
Threads & torch 8; OMP\_NUM\_THREADS unset (BLAS default) \\
Batch & one clip per call, no batching \\
Load average & $0.64$ (1\,min) at start, $1.6$ at end \\
\bottomrule
\end{tabular}
\end{table}

\textbf{What the idle measurement says.} Five of the seven methods run
hundreds to thousands of times faster than real time on this workstation;
MSSA runs at $14\times$ real time and EVMD, at $0.89\times$, is the only
method that cannot keep up with a live 15\,s clip. The wall-clock ranking
is now bandpass $<$ classifier $<$ Conv-TasNet-lite $\approx$ supervised
NMF $<$ standard NMF $<$ MSSA $<$ EVMD, and it is stable: IQRs are under
$1\%$ of the median for every method but supervised NMF ($13\%$) and
MSSA ($4\%$).

\textbf{MACs still do not predict latency, but the example has changed.}
Conv-TasNet-lite has $820\times$ the bandpass filter's MACs and runs
$4.6\times$ slower, not $820\times$ --- because it runs on the GPU. EVMD has
$2.7\times$ Conv-TasNet-lite's MACs and runs $1{,}000\times$ slower ---
because it is per-mode ADMM in NumPy on the CPU. Within one
implementation family the two measures do agree: Baselines~2 and~3, both
this project's NMF code, have a CPU-time ratio of $367/107 \approx
3.4\times$ against a MACs ratio of $5.3\times$. The divergence
\emph{across} implementations is an implementation-efficiency and
hardware effect, which is why Section~\ref{sec:bg-compute} insisted on
measuring both quantities; the loaded run had exaggerated it into
Conv-TasNet-lite being the fastest method of all, which it is not.

\textbf{Scope of these numbers.} Every figure in this subsection is a
desktop measurement: this workstation, running the Python implementations,
idle. They rank the methods against each other on one machine; they are
not a claim about what any of these methods costs on embedded hardware,
where both the processor and the implementation differ. For Baseline~1
that claim now exists separately and is measured rather than extrapolated
--- see Section~\ref{sec:firmware-results}, which supersedes this
subsection's figure for that method alone.

\subsection{The SDR-vs-compute plane}
\label{sec:compute-plane}

Bringing Section~\ref{sec:results}'s separation-quality numbers together
with this section's compute-cost numbers lets every baseline be placed on
a single plane: separation quality (heart SDR, native additive subset) on
one axis, compute cost --- MACs, or measured latency --- on the other.
Two panels are used rather than one collapsed ``compute'' axis, precisely
because Section~\ref{sec:latency} already found MACs and latency
disagreeing substantially for at least one method; collapsing to a single
axis would hide exactly the kind of divergence just demonstrated.

Section~\ref{sec:bg-compute}'s Pareto-dominance definition is applied
directly: on a given axis, baseline A is \emph{dominated} by baseline B if
B has equal-or-better SDR \emph{and} equal-or-lower compute cost, with at
least one strict improvement --- meaning A is never the right choice on
that axis, regardless of how quality and cost are weighted against each
other. Figure/Table generation for this plane (\texttt{sdr\_compute\_plane.py})
reads its SDR from the canonical row-level file
(Table~\ref{tab:baselinevalid}, pooled mean and 95\% CI) and its latency
from the idle-machine run of Table~\ref{tab:latency}, so the plane is
derived from the same artifacts as every other table rather than from
numbers copied into the script.

\begin{table}[H]
\centering
\caption{Pareto dominance per baseline, on the MACs axis and the latency axis separately.}
\label{tab:pareto}
\begin{tabular}{@{}lccc@{}}
\toprule
Method & Heart SDR (dB) & Dominated on MACs? & Dominated on latency? \\
\midrule
Baseline 1 (bandpass) & $5.46 \pm 2.66$ & No & No \\
Baseline 2 (supervised NMF) & $2.58 \pm 1.87$ & Yes & Yes \\
Baseline 3 (standard NMF) & $3.35 \pm 1.88$ & Yes & Yes \\
Baseline 4 (MSSA) & $5.02 \pm 2.91$ & Yes & Yes \\
Baseline 5 (EVMD) & $3.61 \pm 2.35$ & Yes & Yes \\
Baseline 6 (Conv-TasNet-lite) & $5.02 \pm 2.77$ & Yes & Yes \\
\bottomrule
\end{tabular}
\end{table}

\textbf{On the MACs plane, one method dominates outright.} Bandpass
filtering has both the highest heart SDR of the six baselines and, by a
margin of nearly three orders of magnitude, the lowest MAC count
(Table~\ref{tab:macs}) --- every other baseline is strictly worse on both
axes simultaneously, and is therefore Pareto-dominated. If MACs were the
only compute-cost figure available, the conclusion would be unambiguous:
there is no reason, on this evidence, to prefer any learned or iterative
separation method over the zero-training filter.

\textbf{On the measured-latency plane the conclusion is the same, and
an earlier revision of this report said otherwise.} With the loaded-machine
latencies, Conv-TasNet-lite ($23.9$\,ms) appeared faster than bandpass
($36.0$\,ms) and the two were reported as jointly non-dominated --- a
genuine trade-off. On the idle machine bandpass runs in $3.6$\,ms and
Conv-TasNet-lite in $16.7$\,ms on a GPU, so bandpass has the higher SDR
\emph{and} the lower latency and every other baseline is dominated on
both planes. The retracted finding was an artifact of measuring
CPU-bound methods on a saturated CPU while the one GPU-bound method was
unaffected; it is recorded here because it was published in an earlier
revision and because it is a clean example of a plausible measurement
that reversed a conclusion. What remains true is the weaker statement:
Conv-TasNet-lite is the only method whose latency is within an order of
magnitude of the filter's despite three orders of magnitude more
arithmetic, which is what a GPU buys.

% ============================================================
\section{Embedded Deployment: Baseline 1 on a Microcontroller}
\label{sec:firmware}
% ============================================================

Every number in Section~\ref{sec:compute} is a desktop measurement or an
analytic count. Neither answers the question a deployment decision actually
asks: does this run on the hardware it is meant to run on? Section~\ref{sec:related}'s
Edge-Enabled paper \cite{edgelung} targets a PYNQ-ZU FPGA at roughly 15~W;
this section reports Baseline~1 ported to, and measured on, a considerably
smaller target --- an Arduino Nano 33 BLE Sense (nRF52840, 64~MHz Cortex-M4F,
256~kB RAM). The port lives in \texttt{firmware/}, with an ESP32-S3 target
built from the same core but not yet exercised on silicon.

Baseline~1 is the natural first candidate, and not only because it is the
simplest: Section~\ref{sec:compute-plane} puts it alone on the Pareto front
of both compute planes, so it is the method a deployment decision would
consider first.

\subsection{The one intentional divergence: causal filtering}
\label{sec:firmware-causal}

The Python baseline calls \texttt{scipy.signal.sosfiltfilt}, which runs the
cascade forwards and then backwards over the reversed signal. That cancels
the filter's phase response exactly, but it cannot begin until the
\emph{entire} recording exists. The firmware runs the equivalent of
\texttt{sosfilt} instead: one causal pass, block-at-a-time, filter state
carried across calls.

This is a forced choice, not a preference. A 15~s clip at 4000~Hz is 60{,}000
samples; as raw \texttt{int16} that is 120~kB, or 46\% of the Nano 33's total
RAM before any processing begins, and \texttt{filtfilt} additionally needs a
\texttt{float32} working copy of the same signal --- 240~kB more, on a part
with 256~kB in total. It does not fit. Even where it would fit, it forfeits
the point: no output could be produced until the recording ended, which is
buffered batch processing wearing a microcontroller.

Two costs follow, both real. The output carries the filter's phase response
and a non-zero group delay, so waveform-shape comparisons against the Python
baseline are not like-for-like. And \texttt{sosfiltfilt} applies the same
cascade twice, so its stopbands are twice as deep in dB --- one pass of the
order-4 design is a genuinely weaker filter than what Section~\ref{sec:baseline1}
reports.

\textbf{Measured, not assumed.} Running both filters over the same 145
\texttt{Mix.csv} rows with identical coefficients and identical BSS Eval
metrics --- changing only \texttt{sosfiltfilt} to \texttt{sosfilt} --- gives
the mean differences in Table~\ref{tab:causal}.

\begin{table}[H]
\centering
\caption{Causal minus zero-phase, mean over 145 rows. Only the filter direction changes.}
\label{tab:causal}
\begin{tabular}{@{}lccc@{}}
\toprule
Source & $\Delta$SDR & $\Delta$SIR & $\Delta$SAR \\
\midrule
Heart & $+0.18$\,dB & $-0.13$\,dB & $+5.47$\,dB \\
Lung  & $-0.04$\,dB & $-0.50$\,dB & $+7.30$\,dB \\
\bottomrule
\end{tabular}
\end{table}

SDR is unchanged within $\pm0.2$\,dB --- on this dataset the causal filter
costs essentially nothing on the headline metric, better than the theory
above would suggest. SIR drops slightly, exactly as the halved effective
stopband depth predicts, since rejecting the other source is what stopband
depth buys. \textbf{The large SAR gain is a metric artifact and must not be
reported as the firmware separating more cleanly}: \texttt{mir\_eval}'s
projection fits a \emph{causal} 512-tap distortion filter, so a causal IIR's
phase response lies inside the allowed-distortion subspace and is absorbed,
while a zero-phase filter's acausal response does not fit and is partly
scored as artifacts.

\subsection{Verifying a port that has no reference implementation on the target}
\label{sec:firmware-verify}

A build that links is not a build that computes the right thing, and the
board runs a different kernel from the host: CMSIS-DSP's
\texttt{arm\_biquad\_cascade\_df2T\_f32} on the Nano 33, ESP-DSP's
\texttt{dsps\_biquad\_f32} on the ESP32-S3, against portable C on the
desktop. The three sit behind one API so the same source builds for all
three, but the vendor kernels are third-party code, and CMSIS writes the
recursion with the opposite sign convention for $a_1$/$a_2$ from scipy's
--- a detail the initialisation has to correct for and which no compiler
would catch.

The port is therefore verified in four stages, three of which need no
hardware at all:

\begin{enumerate}
  \item \textbf{Host, portable C vs scipy.} The C core filters the same
    golden vectors scipy produced and the deviation is compared against the
    \texttt{float32} quantisation floor. Also checked here: identical output
    for block sizes 1, 7, 64 and 333, which is what proves the streaming
    state is actually carried across calls rather than silently reset.
  \item \textbf{Host, wire protocol.} An emulated board over a pseudo-terminal
    exercises the serial framing in both directions, confirming a faithful
    stream is accepted and a corrupted one rejected rather than silently
    scored.
  \item \textbf{On silicon, vendor kernel.} The \texttt{*-selftest} firmware
    replays a 1024-sample reference held in flash through the vendor kernel
    and prints the deviation over serial. This must be run \emph{before} any
    measurement: if the kernel is wrong, every subsequent number is wrong in
    a way that looks like a filter-design problem.
  \item \textbf{On silicon, full clips.} The \texttt{*-bench} firmware
    replays real HLS-CMDS mixtures compiled into flash and streams the
    filtered output back for scoring.
\end{enumerate}

The bench firmware deliberately has \emph{no microphone}. It filters the same
\texttt{int16} samples the Python baseline filters --- bit-identical, since
the $1/32768$ scaling reproduces \texttt{librosa.load()} exactly --- so the
board's output is diffable against scipy sample by sample and scoreable with
the same BSS Eval metrics as the rest of this report. A microphone would
introduce a transducer, a resampling step and a gain to calibrate away, none
of which the question being asked here needs. Microphone firmwares (PDM on
the Nano 33, I2S on the ESP32-S3) exist for a future real device but are not
in the measurement path.

Only additive mixtures are embedded. Section~\ref{sec:validity} found 109 of
145 rows whose mixture is not $a\cdot(\text{heart}+\text{lung})$ of its named
sources, and an SDR scored against those references measures the dataset's
defect rather than the separator. The first hardware run embedded
\texttt{M0001}--\texttt{M0003}, all three non-additive, which left its SDR
column unquotable; selection now draws from the 36 valid rows.

\subsection{Results on silicon}
\label{sec:firmware-results}

Table~\ref{tab:ondevice} reports the first hardware run: three additive
mixtures replayed from flash on a Nano 33 BLE Sense.

\begin{table}[H]
\centering
\caption{Nano 33 BLE Sense, three additive mixtures replayed from flash.}
\label{tab:ondevice}
\begin{tabular}{@{}lccccc@{}}
\toprule
 & \multicolumn{2}{c}{$|$board $-$ scipy$|$, $\times$ float32 floor} & \multicolumn{2}{c}{SDR (dB), board $=$ scipy} & \\
Clip & heart & lung & heart & lung & \textmu s/sample \\
\midrule
M0087 & 2.04 & 0.89 & 4.918 & 3.686 & 3.474 \\
M0111 & 1.26 & 0.90 & 18.91 & $-12.43$ & 3.488 \\
M0112 & 0.82 & 0.71 & 4.556 & 1.197 & 3.431 \\
\bottomrule
\end{tabular}
\end{table}

\textbf{The port is faithful.} Every SDR agrees with the scipy reference to
the last printed digit, on all six clip/source measurements. These are not
two independent estimates of the same quantity that happen to agree --- they
are the same computation performed on two different machines, and the
agreement is the check.

\textbf{Timing: 3.47\,\textmu s per sample, 72$\times$ faster than real
time.} A 15\,s clip filters in 209\,ms of compute on a 64~MHz Cortex-M4F,
measured around the \texttt{int16} scaling and both band cascades with serial
I/O excluded. That headroom is the deployment-relevant result: this band
split can run as a live stream on a part costing a few dollars, with two
orders of magnitude to spare. For Baseline~1 this figure supersedes
Section~\ref{sec:latency}'s desktop measurement, which timed a different
machine doing a different amount of work. Nothing in the firmware's RAM usage
scales with recording length --- the filters are streaming and the clips live
in flash --- so a 15-second clip and an hour-long stream cost the same few
hundred bytes of filter state.

\subsection{A plausible measurement that was wrong by 18$\times$}
\label{sec:firmware-fpu}

That timing figure was originally 62.4\,\textmu s per sample --- 4.0$\times$
real time. Nothing about it looked wrong. It was a sensible number for a
64~MHz part, it was consistent across all three clips to three significant
figures, and it came from a firmware whose numerical output had already been
verified against scipy twice. It was also 18$\times$ too slow, and the reason
is worth recording because nothing in the measurement itself could have
revealed it.

The discrepancy surfaced only from the resource model of
Section~\ref{sec:firmware-others}, which predicted roughly 400$\times$ real
time for this method. A model being optimistic by 5$\times$ is ordinary; being
optimistic by 100$\times$ is a claim that something is wrong somewhere, and the
model was specific enough to be worth trusting over the intuition that
62\,\textmu s ``seemed fine''.

Reading the actual compiler invocation --- rather than the build
configuration, which is where the intent lives but not the outcome --- showed
the cause. The Arduino-mbed builder appends \texttt{-nostdlib
-mfloat-abi=soft} \emph{after} the framework's own earlier
\texttt{-mfloat-abi=softfp}, and GCC honours the last occurrence of that
flag. Every floating-point operation in the firmware was therefore compiling
to \texttt{\_\_aeabi\_*} software-float library calls. Disassembling both the
project's own filter object and CMSIS-DSP's biquad kernel confirmed it
directly: zero VFP instructions in either, on a Cortex-M4\emph{F} whose
floating-point unit is the entire reason for the suffix.

Adding the correct flag to the build configuration did not fix it, because
those flags are emitted \emph{before} the builder's trailing group and so
still lost the last-occurrence contest --- a detail visible only by
re-reading the compile line after the attempted fix, and easy to miss by
assuming an override had taken effect. PlatformIO's \texttt{build\_unflags}
removes the offending flag outright. After the change: VFP instructions
present in both objects, no software-float calls, 1{,}248 fewer bytes of
flash, and the re-measurement in Table~\ref{tab:ondevice}.

Two things generalise from this. First, a performance figure that looks
plausible has not been validated by looking plausible; here the only thing
that caught it was an independent model disagreeing loudly enough to be worth
chasing. Second, the self-test of Section~\ref{sec:firmware-verify} passed
before and after --- correctly, since software floating point computes the
right answer, only slowly. Numerical verification and performance
verification are separate obligations, and passing the first says nothing
about the second.

\subsection{A tolerance that was calibrated on the wrong signal}
\label{sec:firmware-tolerance}

The first run on additive clips \emph{failed} its accuracy check, at
$1.045\times10^{-5}$ against an absolute tolerance of $10^{-5}$ --- on a
board whose on-silicon self-test had just passed at $5.3\times10^{-8}$ and
whose SDR agreed with scipy to four significant figures. The tolerance was
wrong, not the kernel, and the reason is worth recording because it is a
general trap.

The board computes in \texttt{float32} while the reference is
\texttt{float64}, so the gap between them is not a constant: it scales with
signal amplitude, because an IIR accumulates rounding through its own
recursive state. The $10^{-5}$ figure had been calibrated on
\texttt{M0001}--\texttt{M0003}, which peak near $0.05$. The 36 additive rows
are peak-normalised to full scale --- that normalisation is precisely the
clipping fingerprint that identifies them (Section~\ref{sec:validity}) ---
so they are roughly $20\times$ louder, and the same correct arithmetic lands
roughly $20\times$ further from the reference.

The decisive check is to run \emph{scipy itself} in \texttt{float32} on the
same clip: on \texttt{M0112} that deviates by $1.114\times10^{-5}$, which is
\emph{larger} than the board's $1.045\times10^{-5}$. The board was not merely
acceptable; it was more accurate than the reference implementation would have
been at the same precision.

The criterion is now measured rather than assumed: the \texttt{float32} floor
is computed per clip by filtering the same samples with scipy in
\texttt{float32}, and the board is held to $4\times$ whatever that produces.
Observed ratios on real hardware are $0.71$--$2.04$ (Table~\ref{tab:ondevice}).
For calibration, deliberately injected faults score far higher: a
$0.01\%$ gain error gives $14.3$, a single dropped sample $15{,}691$, and a
sign inversion $296{,}503$ --- so the threshold sits in a wide gap, with
roughly $7\times$ of margin above the real board and two further orders of
magnitude below the smallest fault tested.

\subsection{Server against microcontroller, on identical work}
\label{sec:firmware-vs-server}

The natural question once a method runs in both places is what the move
costs. Answering it needs more care than dividing Section~\ref{sec:latency}'s
desktop figure by the board's, for two reasons that both flatter the board if
left alone.

\textbf{The two platforms were not doing the same work.}
Section~\ref{sec:latency} times \texttt{bandpass\_separate}, which calls
\texttt{sosfiltfilt}: forwards then backwards, two passes per band. The
firmware makes one causal pass (Section~\ref{sec:firmware-causal}). The
desktop figure was therefore for roughly twice the arithmetic, and
\texttt{server\_vs\_board.py} re-times the workstation doing exactly what the
board does rather than halving the existing number and assuming the constant
factors cancel. Measured directly, the zero-phase version costs
$4.8\times$ the causal one, not the $2\times$ its two passes alone would
suggest --- \texttt{sosfiltfilt} also pads the signal, reverses it twice, and
solves for initial conditions per section. Halving the existing figure would
therefore have been wrong by more than a factor of two.

\textbf{A raw ratio conflates two different claims.} 3.6~GHz against 64~MHz
is a $56\times$ clock disadvantage before architecture enters the picture.
``This processor is faster'' and ``this processor does more per cycle'' are
separate statements, and only the second says anything about whether the port
was a good idea. Cycles per sample separates them.

\begin{figure}[H]
\centering
\includegraphics[width=\textwidth]{../results/plots/server_vs_board.pdf}
\caption{The $276\times$ throughput gap decomposed. Because the two factors
multiply rather than add, the top panel places them on a log axis, where the
segment lengths are exactly $\log 56$ and $\log 4.9$. The lower panel is the
measurement underneath: the same causal band split, normalised for clock rate,
where the gap is a factor of five rather than a factor of 276.}
\label{fig:servervsboard}
\end{figure}

\begin{table}[H]
\centering
\caption{Identical workload --- one causal pass per band over the same 60{,}000-sample mixture. Server figures are the median of five independent runs (spread under 1\%) at load average 0.1; board figures measured on silicon.}
\label{tab:servervsboard}
\begin{tabular}{@{}lrrrr@{}}
\toprule
Platform & ms/clip & \textmu s/sample & cycles/sample & $\times$ real-time \\
\midrule
Server, i9-9900K @ 3.6\,GHz        & 0.75 & 0.0126 & 45.2 & 19{,}900 \\
Board, Cortex-M4F @ 64\,MHz        & 208.5 & 3.474 & 222.3 & 72 \\
\addlinespace
Board, before the FPU fix          & 3{,}747 & 62.45 & 3{,}997 & 4.0 \\
\bottomrule
\end{tabular}
\end{table}

\textbf{Most of the gap is clock rate, but not all of it.} The board is
$275\times$ slower per sample against a $56\times$ clock disadvantage, so per
clock cycle it spends $4.9\times$ what the workstation does on the same
filtering. Decomposing the $275\times$ that way is the useful move: roughly a
fifth of it is architecture and the rest is simply being clocked slower.

Both parts are explicable at the level of the inner loop. Each output sample
needs eight biquad sections (four per band, two bands), so the workstation
spends about $5.7$ cycles per section and the board about $28$. A
direct-form-II transposed biquad is five multiply-accumulates with a serial
dependency between consecutive samples: the $5.7$ is close to what a
dependency-limited chain costs on an out-of-order core, and the $28$ is the
expected cost of the same five operations as VFP instructions on an in-order
Cortex-M4F, plus loop overhead. Neither number is anomalous, which is the
point --- the comparison is measuring the two architectures rather than an
artefact.

This is the substantive argument for the port, stated at its real strength: a
microcontroller gives up about a factor of five in per-cycle efficiency here,
not the factor of $275$ the raw timing suggests, and $72\times$ real-time
headroom makes that affordable. The serial structure of an IIR cascade is
what keeps the architectural penalty this small --- issue width and
out-of-order execution have a dependency chain to fight. A method with
exploitable instruction-level parallelism, or one whose working set fits a
desktop cache but not 256~kB of SRAM, would show a much worse ratio;
Section~\ref{sec:firmware-others} covers what happens to those methods here
for a more basic reason.

The last row is the same silicon before the build fix of
Section~\ref{sec:firmware-fpu}, at $3{,}997$ cycles per sample --- an
$88\times$ per-cycle penalty against the workstation rather than
$4.9\times$. It is kept because it is what this comparison would have
concluded from a measurement that looked entirely reasonable at the time: not
``the toolchain is misconfigured'' but ``microcontrollers are getting on for
two orders of magnitude less efficient per cycle than desktops, which is
simply what embedded DSP costs''. That conclusion would have been wrong,
quotable, and very hard to dislodge later.

The server side of this table needed the same correction for the same
reason. Its first version timed a function that designed the filter
coefficients inside the timed call --- work the firmware never does, since
its coefficients are compile-time constants in flash --- and so charged the
workstation $2.8\times$ its real cost, which would have reported the
architectural penalty as $1.7\times$ rather than $4.9\times$. That figure was
also stable across repetitions and plausible on its face. Two independent
instances in one comparison is a reasonable argument for treating ``the
number looks sensible'' as carrying no evidential weight at all.

\subsection{Energy, where the answer is not yet known}
\label{sec:firmware-energy}

Throughput is the wrong axis for a battery-powered device, which is judged on
joules. The two need not agree: a processor $276\times$ slower that draws
$1000\times$ less power still wins on energy. Half of that comparison can be
measured here and half cannot, and the half that cannot turns out to be
decisive.

\textbf{The server side is measured, with no instrumentation.} Intel CPUs
expose a package energy counter through RAPL as an ordinary file in
\texttt{/sys}, in microjoules, so sampling it either side of a known workload
gives real energy rather than a TDP estimate. Idle package power is
$28.0$\,W; under the board's exact workload it is $46.7$\,W, so the
computation itself draws $18.7$\,W above idle. At $0.75$\,ms per clip that is
$\mathbf{14.3}$\,\textbf{mJ} marginal, or $35.8$\,mJ counting the whole
package --- median of five runs, spread $2\%$. Both are reported because the
right one depends on the question: marginal if the server exists anyway and
one more clip is the question, full package if a dedicated device is being
compared against a dedicated machine.

\textbf{The board side is not measured, and cannot be from here.} The Nano 33
BLE Sense has no onboard current sense, so the firmware cannot report its own
consumption and nothing on the server can infer it. What is reported instead
is the \emph{break-even}: since the board's time per clip is known precisely
($208.5$\,ms), energy is that time multiplied by the one missing number, and
parity arrives at $\mathbf{69}$\,\textbf{mW} against the server's marginal
cost, or $172$\,mW against its full package.

\textbf{That number is why no datasheet figure is substituted.} A development
board with USB, regulators, onboard sensors and an LED awake can plausibly
sit on either side of $69$\,mW, so the conclusion flips inside the range a
reasonable assumption would cover. Where a comparison is insensitive to an
assumption, assuming is fine; this one inverts, and an inverted conclusion
asserted from a datasheet would be worse than an acknowledged gap. An inline
USB power meter during a \texttt{*-bench} run would settle it to the
order of magnitude the argument needs.

Two things are worth noting about what such a measurement would show. If the
board does come in under $69$\,mW, then the platform that is $276\times$
slower is simultaneously cheaper per clip, and the deployment argument
reverses relative to Table~\ref{tab:servervsboard} --- which is the result
this section exists to flag as plausible rather than to claim. And the same
measurement would price the build flag of Section~\ref{sec:firmware-fpu}: the
pre-fix binary ran $18\times$ longer for identical output, so at comparable
power draw it also cost $18\times$ the energy, which on a battery is the
difference between a day and a fortnight.

\subsection{Why the other five baselines do not run here}
\label{sec:firmware-others}

Porting Baseline~1 raises the obvious question: is there anything
\emph{better} that would also fit? Attempting five more ports to find out
would be expensive, so \texttt{src/mcu\_feasibility.py} models each method's
peak working memory, parameter storage and compute time against both targets,
deriving every figure from the method's own imported constants.

\textbf{The discriminator is not arithmetic.} A method is \emph{streaming} if
it produces output from a bounded window of input while carrying only fixed
state between calls, and \emph{batch} if the whole recording must be resident
before any output exists. Baseline~1 is streaming: 16 floats of biquad state,
identical whether the recording is 15 seconds or an hour. That property, not
the MAC count, decides most of Table~\ref{tab:mcufeas}.

\begin{table}[H]
\centering
\caption{Peak working RAM against the 211~kB left on a Nano 33 BLE Sense after the Baseline 1 firmware's own footprint. All figures at \texttt{float32}.}
\label{tab:mcufeas}
\begin{tabular}{@{}lccl@{}}
\toprule
Method & Streaming & Peak RAM & Verdict \\
\midrule
Baseline 1 (bandpass)        & yes & 832\,B  & fits ($0.004\times$) \\
Baseline 2 (supervised NMF)  & yes & 9.1\,kB & fits ($0.04\times$) \\
Baseline 3 (standard NMF)    & no  & 1.0\,MB & $4.7\times$ over \\
Baseline 4 (MSSA)            & no  & 36\,MB  & $167\times$ over \\
Baseline 5 (EVMD)            & no  & 15.4\,MB & $71\times$ over \\
Baseline 6 (Conv-TasNet-lite) & no & 9.2\,MB & $43\times$ over \\
\bottomrule
\end{tabular}
\end{table}

\textbf{Baselines 4 and 5 are batch by construction, not by implementation.}
SSA embeds the signal into an $L \times K$ trajectory matrix ($L{=}50$,
$K{=}59{,}951$), takes its SVD, and reconstructs each component by diagonal
averaging back across the full length; every one of those structures is
$O(LN)$, and the reconstruction references the whole signal. VMD's ADMM
iterates in the frequency domain over the mirror-padded signal, holding all
$K$ modes and their spectra live across every iteration. Neither has a
streaming formulation to be written more carefully --- the algorithms
themselves are defined over the complete recording.

\textbf{Baseline 6 is batch only as written.} It encodes the whole clip and
runs the TCN over all $6{,}000$ resulting frames at once, so its $9.2$\,MB is
activations rather than weights --- the parameters are $318$\,kB at
\texttt{int8} and fit flash comfortably. A chunked, causal Conv-TasNet is a
known variant and would make activations $O(\text{chunk})$; that is a
redesign of the model, not a deployment step, and it is the honest description
of what stands between this method and the board.

\textbf{Quantisation does not rescue any of them.} Moving from
\texttt{float32} to \texttt{int8} divides memory by four, which leaves
Baseline~3 at $1.2\times$ over, Baseline~6 at $10.7\times$, Baseline~5 at
$17.8\times$ and Baseline~4 at $41.6\times$. None fits. More fundamentally,
for most of them \texttt{int8} is not an applicable representation at all:
Baseline~5's ADMM convergence tolerance is $10^{-6}$, four orders of magnitude
below \texttt{int8}'s resolution of roughly $4 \times 10^{-3}$, so the solver
could not detect its own convergence; Baseline~4 selects components by a
$2\%$ eigenvalue-contribution threshold across singular values spanning many
orders of magnitude, which 256 levels cannot represent; and the NMF
multiplicative updates carry a $10^{-12}$ epsilon floor. Only Baseline~6 is a
natural quantisation target --- neural post-training quantisation being
mature and well tooled --- and for it quantisation is the \emph{last} step
after the streaming redesign, not an alternative to it.

\textbf{What this means for the deployment question.} On the native additive
subset (Table~\ref{tab:sixmethodnative}), the only method that beats
Baseline~1 on either source is Baseline~4 on lung, by $1.34$\,dB --- and
Baseline~4 is the least feasible entry in the table. Every method that could
run here scores worse than Baseline~1 on both sources. The finding is
therefore a negative one, and worth stating plainly: on this dataset, with
these six methods, the cheapest method is also the best-performing one, which
is why Section~\ref{sec:compute-plane} already placed it on the Pareto front
of both compute planes. Baseline~2 remains the one non-trivial port that
would fit, and its value would be demonstrating that the firmware's
verification chain generalises to a fundamentally different algorithm ---
FFT and iterative matrix updates rather than a biquad cascade --- not
improving separation quality.

\textbf{One caveat on the timing column.} The model's compute estimates
assume one MAC per cycle with operands in registers. Calibrated against the
one method now measured on silicon, it is optimistic by $5.6\times$
($400\times$ predicted, $72\times$ measured). Applying that factor to
Baseline~2 would put it at $0.8\times$ real time rather than $4.4\times$ ---
so its real-time margin is genuinely uncertain, and would need measuring
rather than asserting. The memory columns, which decide every infeasible row,
are unaffected.

% ============================================================
\section{Related Work}
\label{sec:related}
% ============================================================

Table~\ref{tab:related} summarizes six papers reviewed against this work,
organized by what each one measures. Full citation details for all six are
now filled in below, sourced from the Notion Reading List; five of the six
have since been independently cross-checked against their primary-source
PDFs rather than relied on second-hand summaries (\cite{aidriven}, a
Ph.D.\ dissertation, remains the one exception --- flagged in the
bibliography).

\begin{table}[H]
\centering
\caption{Related work, organized by what is measured.}
\label{tab:related}
\resizebox{\textwidth}{!}{%
\begin{tabular}{@{}p{3.3cm}p{2.1cm}p{2.6cm}p{2.6cm}p{2.6cm}p{3cm}@{}}
\toprule
Paper & Dataset & Separation method & Separation metrics & Downstream task & Split \\
\midrule
HLS-CMDS descriptor \cite{hlscmds} & --- (source) & --- & --- & --- & --- \\
AI-Driven Cardiorespiratory Signal Processing \cite{aidriven} & HLS-CMDS & PL-NMF + LLM-guided fundamental frequency (``LingoNMF'') & not reported (paired $t$-test on unspecified proxy) & clustering + anomaly detection, not accuracy & not stated \\
Spectrotemporal DL for Heart Sound Classification \cite{spectrotemporal} & PhysioNet 2016 (training) + \textbf{HLS-CMDS} (external validation, Sec.~5 of the paper) & \textbf{bandpass filter} (their Experiment 4 stress test: heart sounds computationally separated from mixed HLS-CMDS heart-lung recordings) & none (no SDR/SIR/SAR --- only downstream accuracy is reported) & 4-class PCG classification (Normal/Murmur/Extra~Sound/Rhythm~Disorder); \textbf{accuracy collapses 89.0\%$\to$41.0\% (Exp.\,3 clean vs.\ Exp.\,4 separated) --- see Sec.~\ref{sec:yaqub} below, this is the motivating result for this project's whole research question} & not stated \\
NMF-Enhanced Log-Mel + CRNN \cite{nmfcrnn} & ICBHI 2017 + Fraiwan CWLS & NMF-based enhancement & none (no SNR/SDR/correlation) & multi-class respiratory disease classification & recording-level (leakage flagged by authors) \\
Edge-Enabled Portable Lung-Sound Classifier \cite{edgelung} & \textbf{HLS-CMDS} & proposed lung isolation from mixtures & \textbf{none} & lung sound classification & not stated \\
Performance Measurement in Blind ASS \cite{vincent2006} & --- (methods) & --- & defines SDR/SIR/SAR & n/a & n/a \\
Cardiorespiratory Separation via SSA \cite{ssa} & not stated & two-stage Singular Spectrum Analysis & SDR, STOI, correlation vs.\ 5 baselines & none & n/a \\
\bottomrule
\end{tabular}%
}
\end{table}

\subsection{Yaqub et al.: the motivating result, not an unrelated
classification paper}
\label{sec:yaqub}

An earlier draft of this table filed \cite{spectrotemporal} under
``PhysioNet 2016, no separation, binary normal/abnormal classification.''
That was wrong on every count that matters, confirmed now by reading the
primary source directly rather than a second-hand summary. Yaqub et al.
train a ResNet-18 PCG classifier on PhysioNet 2016, then use HLS-CMDS
purely for \emph{external validation} across four experiments: (1) binary
fine-tuning, 88.0\% accuracy; (2) 4-class fine-tuning, 86.0\%; (3) full
retraining from ImageNet weights on HLS-CMDS's own 4-class labels, 89.0\%
(their best, most balanced result); and (4) a stress test that re-evaluates
the Experiment-3 model on heart sounds \emph{computationally separated from
HLS-CMDS's mixed heart-lung recordings using a standard bandpass filter} ---
where accuracy collapses to 41.0\% (Table~9 of \cite{spectrotemporal}: n=5365,
weighted F1 falls from 0.89 to 0.39, every class's precision and recall
degrade, not just one).

This is not a tangential result. It is the single empirical finding this
project's entire research question exists to explain: separating a mixed
cardiorespiratory recording before classifying it costs something, and
Yaqub et al.\ are the first to measure that cost directly, end to end, on
this project's own dataset. Their own Discussion (Sec.~6.2) attributes the
collapse specifically to ``the subtle distortions and potential information
loss introduced by the standard bandpass filter used for source
separation'' and its ``inherent limitation... to perfectly separate
spectrally overlapping signals'' --- which is exactly the same class of
method as this report's own Baseline~1 (Section~\ref{sec:baseline1}), and
exactly the SIR-vs-SAR tradeoff this project's own bandpass results
demonstrate directly on the separation-quality side (real SIR gain, real
SAR cost). Yaqub et al.\ show what that tradeoff costs a downstream
classifier; this project's separation-quality metrics are the mechanism
that predicts when it will. Their 89\%$\to$41\% is one point on the curve
this project's Condition~A/B comparison (Section~\ref{sec:conclusion}) is
built to trace out continuously against SDR, not a two-point anecdote to
reproduce once.

\subsection{Metrics caveats from the BSS Eval literature}
\label{sec:relatedmetrics}

Two points from \cite{vincent2006} directly shaped the design in
Section~\ref{sec:metrics}--\ref{sec:baseline1}: (1) SAR is sensitive to
distortion that a downstream classifier's accuracy may not register at all
(a classifier can stay correct despite real signal artifacts, if the
distortion doesn't cross a decision boundary), which motivates reporting
SDR/SIR/SAR rather than accuracy alone; and (2) SDR is not a perceptual
measure and is gameable --- a mildly low-pass-filtered estimate can drive
SDR toward $+\infty$ under a time-varying filter without the estimate
actually improving --- which is why Tables~\ref{tab:baselinefull}
and~\ref{tab:baselinevalid} are reported with all three metrics rather than
SDR in isolation. The same source notes
that HLS-CMDS mixtures are heart+lung only, with no independent
sensor-noise reference, so the noise term in a full images-based BSS Eval
decomposition is effectively zero here; SAR in this report reflects
algorithmic artifact, not resilience to real acquisition noise.

One further cross-check surfaced while reviewing \cite{edgelung}: that
paper's stated HLS-CMDS lung class counts (Normal 28, Wheezing 28, Pleural
Rub 25, Rhonchi 23, Fine Crackles 22, Coarse Crackles 19) match the
\emph{total-in-Mix.zip} counts already encoded in this repository's
\texttt{load\_dataset.py} (its \texttt{PAPER\_LUNG\_COUNTS}
\texttt{mix\_zip} column, validated against the live CSVs), not the smaller
isolated-\texttt{LS.csv} count (5--12 per class). The two counts are easy to
conflate and any report drawing on this dataset should be explicit about
which one it cites.

% ============================================================
\section{Discussion}
\label{sec:discussion}
% ============================================================

Of the six papers reviewed, exactly one --- \cite{edgelung} --- evaluates
on HLS-CMDS \emph{and} classifies mixture-separated audio, but reports zero
separation-quality metrics and, per the summary reviewed, does not compare
against classifying the isolated ground truth directly. The SSA paper
\cite{ssa} reports proper separation metrics but no classification task at
all; the NMF+CRNN paper \cite{nmfcrnn} reports classification accuracy with
real numbers but on a different dataset, with no separation metrics, and its
own authors flag a recording-level (leakage-prone) split as a limitation.
No paper reviewed reports separation quality, separated-audio classification
accuracy, and isolated-ground-truth classification accuracy together, under
a split verified not to leak recordings between the pool used to fit
anything and the pool used to evaluate it.

That fourth condition is not a hypothetical concern for this dataset
specifically: Section~\ref{sec:leakage}'s audit found that a third of the
isolated recordings are literal duplicates of recordings embedded in the
mixtures, and that mix rows chain together into leak groups as large as 32
of 145 rows. A recording-level split --- the exact limitation \cite{nmfcrnn}
flags in its own results --- would leak ground truth into training on this
dataset in a very concrete way.

The Baseline~1 result (Section~\ref{sec:results}) demonstrates the
SIR/artifact trade-off Section~\ref{sec:relatedmetrics} predicts on real
data rather than asserting it from theory: a real, if crude, separation
method's SDR gain over doing nothing is much smaller than its SIR gain alone
would suggest, because SAR moves sharply in the opposite direction. This is
the concrete case for reporting classification accuracy alongside SDR/SIR/SAR
in the planned next phase (see \texttt{PROTOCOL.md}) rather than either
metric family in isolation: a separation method that trades SIR for SAR in
this way may or may not hurt a downstream classifier, and that is an
empirical question this pipeline is built to answer, not one this report
answers on its own.

Baseline~2's comparison to Baseline~1 raises a second, more specific
concern. On the full row set the two methods were close to tied; on the
36-row subset where genuine separation is actually possible
(Section~\ref{sec:validity}), leakage-safe supervised NMF is clearly
\emph{worse} than the zero-training filter, not comparable to it. Two
non-exclusive explanations follow from how this pipeline is built. First,
closing the leakage gap (Section~\ref{sec:leakage}) removes exactly the
kind of inflation a na\"ive reproduction would have kept --- per this
project's own working assumption, a lower, leakage-safe number is the
\emph{correct} one, not evidence the method itself is broken. Second, the
gap is far larger and more consistent on the valid subset than on the full
set, which is the opposite of what pure leakage inflation alone would
predict (leakage should inflate NMF's fitted-dictionary case regardless of
row validity); this instead points toward the fixed dictionary rank
($K_i{=}10$, $K_r{=}20$) or the fixed iteration counts
(Section~\ref{sec:baseline2}) not transferring well to this dataset's real
acoustic mixtures, as opposed to the synthetically-summed or curated
mixtures a published method may have been tuned against. A third, colder explanation was raised informally while Baseline~2 was
being built: an \emph{oracle} mask built from one mixture row's own true
heart/lung spectrograms --- cheating on purpose, purely to bound the
method's ceiling --- appeared to beat the no-separation reference only
modestly, which would have meant the dataset's real mixture recordings
impose a hard ceiling on mask-based separation that no amount of NMF
sophistication fixes. That probe is now a committed script
(\texttt{oracle\_mask\_ceiling.py}) run over all 36 valid rows on
Baseline~2's own STFT grid, and it does not survive the re-check. The ideal
ratio mask --- exactly Baseline~2's mask form with the true magnitudes
substituted --- scores $11.5 \pm 1.9$\,dB heart SDR and $9.2 \pm 1.7$\,dB
lung (95\% CI over rows), against $2.5$ and $-2.0$\,dB for the raw mixture
and $5.5$ and $4.1$\,dB for bandpass; the ideal Wiener filter reaches
$12.4$ and $10.6$\,dB. Baseline~2 therefore sits $7.8$\,dB (heart) and
$9.7$\,dB (lung) below the ceiling of its own method family, and the
shortfall is the dictionaries' and activations', not the dataset's. The
one-row note was almost certainly made on a non-additive row, where
Section~\ref{sec:validity} shows every method --- an oracle included ---
inherits the unrelated mixture's energy as unremovable distortion. The
check does establish a ceiling of a different kind: a magnitude mask on a
512-point STFT tops out near $12$\,dB on these recordings, which is below
the $15$--$20$\,dB region where Section~\ref{sec:knee}'s first classifier
starts to prefer the output of a spectral-stripping separator. Whether a finer time-frequency grid or a
time-domain separator can pass that cap is not tested here.

% ============================================================
\section{Limitations and Future Work}
\label{sec:conclusion}
% ============================================================

\begin{itemize}
  \item \textbf{The descriptor's stated sample rate is unreconciled.} The
    released dataset measures 4000\,Hz throughout
    (Section~\ref{sec:provenance}) against the 22{,}050\,Hz its descriptor
    paper states. Provenance itself is settled --- the copy audited here is
    byte-identical to the Mendeley release --- so the open question is
    narrower than it first appeared: whether the descriptor misstates the
    rate, or whether the figure was mis-transcribed downstream of it.
    Resolving that against the descriptor's own text, with a page reference,
    is outstanding, and it determines how strongly the companion
    dataset-audit comment can phrase the recommendation.
  \item \textbf{The trajectories are complete but their vertical
    resolution is 36 rows, and a crossing is a point estimate from seven
    sampled levels.} Section~\ref{sec:knee}'s trajectories and
    Section~\ref{sec:robustness}'s repeat of them are now full six-baseline
    results, and both say the same thing qualitatively. But every point on
    them is an accuracy over 36 recordings --- one row is $2.8$ points, and
    a binomial 95\% interval near $40\%$ is about $\pm16$ points --- which
    is why the crossing \emph{values} move by $6$--$11$\,dB between two
    classifiers, why their bootstrap intervals reach the top of the
    measured grid (Table~\ref{tab:kneeboot}), and why this report quotes a
    crossing region under a stated interpolation protocol, not a threshold.
    The 36 rows are also not 36 independent units: they fall into 24 leak
    groups (one of 5, one of 3, six of 2, sixteen singletons), so the
    effective sample size lies between 36 (no within-group correlation)
    and $17.5$ (Kish, perfect correlation); the empirical within-group
    correlation of per-row correctness is $0.24$ on isolated audio and
    near zero on separated audio, which puts it near $29$. The
    leak-group bootstrap of Table~\ref{tab:kneeboot} already accounts
    for this and widens nothing by more than a decibel.
    Tightening that needs more valid rows than HLS-CMDS has, which is the
    synthetic substrate's job (Section~\ref{sec:synthetic}); running the
    sweep and both classifiers over its 1500 rows is the direct next step.
  \item \textbf{Architecture 2 was tuned once, by hand.}
    Section~\ref{sec:robustness}'s CNN went from $30.0\%$ to $62.0\%$
    through feature standardisation and a longer mini-batch schedule, with
    three configurations tried and the selection made on training-set fit;
    that is a small, disclosed search, but it is a search, and the adopted
    configuration was also the one with the best test accuracy at the
    time. The level normalisation added afterwards left its accuracy
    unchanged and narrowed its CI to $\pm3.9$ points. A nested selection,
    or simply a second seed, would say how much of the $62.0\%$ is the
    configuration and how much is the fold draw --- neither has been run.
  \item \textbf{The level confound was caught by its symptom, not by
    design.} Section~\ref{sec:conditionb}'s first run was wrong in a way
    that produced a plausible table, and it was found only because three
    unrelated methods returned identical predictions. The classifiers are
    now level-invariant, but nothing yet checks that separated inputs
    fall inside the region where the SVM's kernel is informative; a
    standardised-distance report per condition would make the next such
    defect visible before it is published.
  \item \textbf{Six separation methods evaluated; none dominates cleanly.}
    Baselines~1, 4, and 5 are zero-training sanity floors by design;
    Baselines~2, 3, and 6 involve some form of fitting. Section~\ref{sec:sixmethod}
    shows Baseline~6 (Conv-TasNet-lite) scores best on synthetic-substrate
    SDR for both sources and on native heart SDR, but loses to Baseline~4
    (MSSA) on native lung SDR by more than either method's CI half-width ---
    so this pipeline still does not have a single method that demonstrably
    and consistently improves on the strongest zero-training baseline across
    both sources and both substrates. LLM-guided NMF (``LingoNMF''), the one
    reviewed method not yet reproduced (Table~\ref{tab:related}), remains an
    open candidate.
  \item \textbf{Baseline 6 is an undertrained first neural baseline, not a
    tuned model.} Table~\ref{tab:b6diag} shows every fold exhausts its
    25-epoch budget without the early-stopping criterion ever firing, on a
    pool of $\sim$32 training recordings per class per fold --- consistent
    with a compute-budget-limited result rather than a converged one
    (Section~\ref{sec:baseline6}). Its results should be read as
    infrastructure validation, not as this project's best achievable neural
    result; a longer training budget and/or more aggressive on-the-fly
    augmentation are the natural next steps before treating Baseline~6's
    ranking against Baselines~1--5 as final.
  \item \textbf{Baseline 5's synthetic-substrate column is a subsample, not
    the full set.} The full $K{=}2..10$ VMD sweep at the full synthetic
    $n{\approx}1500$ was measured at $\sim$15\,s/mixture
    (Section~\ref{sec:baseline5}), extrapolating to multi-hour runtime; its
    synthetic column in Table~\ref{tab:sixmethodsynth} therefore runs on a
    stratified 75-row subsample, giving it a visibly wider CI than the
    other five methods' 1500-row columns. This is a disclosed scope
    reduction, not a like-for-like comparison, and should not be read as
    Baseline~5 being noisier in an absolute sense.
  \item \textbf{The synthetic set's gain range is fitted from a small
    sample.} \texttt{synthetic\_mix.py}'s gain distribution
    (Section~\ref{sec:synthetic}) is fitted from the 36 native additive rows
    rather than chosen arbitrarily, which is a strength of the construction
    and equally a dependency of it: the synthetic substrate inherits whatever
    is idiosyncratic about those 36 rows' own mixing gains
    (3.5--135$\times$), and $n{=}36$ is a small sample to fit a range from. A
    sensitivity check --- rebuilding the synthetic set under a deliberately
    wider and a deliberately narrower gain range, then re-running one
    baseline --- would bound how much Table~\ref{tab:sixmethodsynth} depends
    on that fit.
  \item \textbf{Baseline 2 hyperparameters: verified against the primary
    source, with two corrections.} An earlier revision listed the NMF
    dictionary ranks and iteration counts (Section~\ref{sec:baseline2}) as
    transmitted second-hand and unverified. They have since been checked
    against \cite{nmfcrnn} directly (its Sections~3.2.3--3.2.4): $K_i{=}10$,
    $K_r{=}20$, 100 and 60 multiplicative-update iterations and the STFT
    window are confirmed exact, and the check found two discrepancies that
    were corrected --- the STFT hop had been an unverified 128 (the paper's
    is 256), and the paper bandpass-denoises every snippet at 50--1800\,Hz
    before the STFT, a stage the first reproduction lacked. Neither change
    moved Baseline~2's full-set SDR by more than the fold-to-fold spread, so
    the method's shortfall against Baseline~1 was not caused by either.
    The paper's own training-dataset description remains ambiguous between
    ICBHI\,2017+Fraiwan and HLS-CMDS's isolated recordings
    (Table~\ref{tab:related}); this project fits on HLS-CMDS's isolated
    recordings under the leakage-safe pool and says so.
  \item \textbf{The oracle-mask ceiling is now measured, and it refutes
    the informal note that prompted it.} \texttt{oracle\_mask\_ceiling.py}
    runs three oracle masks over all 36 valid rows (Section~\ref{sec:discussion});
    the ideal ratio mask reaches $11.5$\,dB heart SDR against $2.5$\,dB for
    the raw mixture, so there is no dataset-imposed ceiling near the
    no-separation floor. What the check does establish is a different
    ceiling: a magnitude mask on a 512-point STFT cannot exceed roughly
    $12$\,dB here, which is below Architecture~1's knee region. Whether a
    finer grid or a time-domain method lifts that cap is untested.
  \item \textbf{No sensor-noise robustness claim.} As noted in
    Section~\ref{sec:relatedmetrics}, HLS-CMDS mixtures contain no
    independent noise source, so this pipeline's SAR does not speak to
    robustness against real acquisition noise (relevant to
    hardware-deployed classifiers such as \cite{edgelung}'s).
  \item \textbf{Citation completeness.} Full bibliographic details (author,
    venue, year) for four of the six related-work papers were not available
    when this report was drafted (\cite{nmfcrnn} has since been resolved,
    Section~\ref{sec:related}); the remaining entries should be completed
    before any external submission.
  \item \textbf{Table 6 typography.} The related-work table is scaled to
    the page width via \texttt{resizebox} to fit six columns, which produces
    a few awkward line breaks inside narrow cells (e.g.\ ``HLS-CMDS'',
    ``SNR/SDR/correlation''). Cosmetic only, but worth widening manually
    (e.g.\ splitting into two tables, or rotating to landscape) before any
    polished submission.
  \item \textbf{No CI wiring.} The 31-test suite (Section~\ref{sec:validation})
    only runs when invoked manually (\texttt{make test}).
  \item \textbf{Small heart-sound classes threaten the planned classifier's
    per-fold coverage.} Table~\ref{tab:perclass}
    (Section~\ref{sec:perclass}) shows S4 ($n{=}2$), AV Block, and
    Tachycardia ($n{=}3$ each) are far smaller than the other heart classes.
    Under the leak-group-respecting 5-fold split, a 2-recording class cannot
    be represented in every fold's training set; the classification stage
    (\texttt{PROTOCOL.md} \S5.3) should decide up front whether to merge,
    exclude, or specially stratify these classes rather than discover the
    problem as an unexplained per-class accuracy anomaly later.
\end{itemize}

% ============================================================
\begin{thebibliography}{9}

\bibitem{hlscmds}
Y.~Torabi, S.~Shirani, and J.~P.~Reilly, ``Descriptor: Heart and Lung Sounds
Dataset Recorded from a Clinical Manikin using Digital Stethoscope
(HLS-CMDS),'' \textit{IEEE Data Descriptions}, 2025.
doi:10.1109/IEEEDATA.2025.3566012.

\bibitem{vincent2006}
E.~Vincent, R.~Gribonval, and C.~F\'{e}votte, ``Performance measurement in
blind audio source separation,'' \textit{IEEE Transactions on Audio,
Speech, and Language Processing}, vol.~14, no.~4, pp.~1462--1469, 2006.

\bibitem{mireval}
C.~Raffel et al., ``mir\_eval: A Transparent Implementation of Common MIR
Metrics,'' in \textit{Proc.\ ISMIR}, 2014.

\bibitem{luomesgarani2019}
Y.~Luo and N.~Mesgarani, ``Conv-TasNet: Surpassing Ideal Time--Frequency
Magnitude Masking for Speech Separation,'' \textit{IEEE/ACM Transactions on
Audio, Speech, and Language Processing}, vol.~27, no.~8, pp.~1256--1266,
2019.

\bibitem{aidriven}
Y.~Torabi, ``AI-Driven Cardiorespiratory Signal Processing: Separation,
Clustering, and Anomaly Detection,'' Ph.D.\ dissertation, McMaster
University, 2025. (Cited via the Notion Reading List; not yet
independently read from the primary source --- see
Section~\ref{sec:conclusion}.)

\bibitem{spectrotemporal}
A.~Yaqub, M.~S.~Orakzai, M.~F.~Qureshi, Z.~Mushtaq, I.~Siddique, and
T.~Radwan, ``Spectrotemporal Deep Learning for Heart Sound Classification
under Clinical Noise Conditions,'' \textit{Computer Modeling in
Engineering \& Sciences}, vol.~145, no.~2, pp.~2503--2529, 2025.
doi:10.32604/cmes.2025.071571. (Independently read from the primary
source --- see Section~\ref{sec:related}.)

\bibitem{nmfcrnn}
B.~Han, W.~Quan, B.~Matuszewski, and D.~Corbett, ``Respiratory Disease
Classification Using NMF-Enhanced Log-Mel Spectrograms and Convolutional
Recurrent Neural Networks,'' \textit{Sensors}, vol.~26, no.~13, art.~4268,
2026. doi:10.3390/s26134268. (Independently read from the primary source.)

\bibitem{edgelung}
A.~Puneet, P.~Shankar, S.~R.~R.~Koluguri, and A.~Srivastava,
``Edge-Enabled Portable Classifier for Lung Sounds Using Convolutional
Neural Networks,'' in \textit{Proc.\ 2025 IEEE Biomedical Circuits and
Systems Conference (BioCAS)}, pp.~21--25, 2025.
doi:10.1109/BioCAS67066.2025.00016. (Independently read from the primary
source.)

\bibitem{ssa}
B.~Han and W.~Quan, ``Cardiorespiratory Sound Separation Using Singular
Spectrum Analysis,'' in \textit{Proc.\ 2025 17th Int'l Conf.\ on Signal
Processing Systems (ICSPS)}, 2025. doi:10.1109/ICSPS66615.2025.11347745.
(Independently read from the primary source.)

\end{thebibliography}

\end{document}
